% Volume unificado do Manual de TCC da UnDF e do Guia de Uso do UnDF-TeX.
% Os documentos autônomos permanecem em manuais/parte-i e manuais/parte-ii.

\DocumentMetadata{lang=pt-BR,tagging=on}
\tagpdfsetup{math/mathml/luamml/load=false}
\documentclass[12pt,a4paper,oneside]{report}
\usepackage{estilos/manual-style-integrado}
\usepackage{pdflscape}
\usepackage{amssymb}
\usepackage{fancyvrb}

\hypersetup{
  pdftitle={Manual de Trabalhos de Conclusão de Curso e Guia de Uso do UnDF-TeX},
  pdfauthor={Universidade do Distrito Federal Professor Jorge Amaury Maia Nunes},
  pdfsubject={Orientações acadêmicas, metodológicas, normativas e técnicas},
  pdfkeywords={UnDF, UnDF-TeX, TCC, trabalhos acadêmicos, ABNT, versão 1.0.0}
}

\titleformat{\part}[display]
  {\centering\normalfont\bfseries\Large}
  {PARTE \thepart}
  {0.8em}
  {\MakeUppercase}
\titlespacing*{\part}{0pt}{0.20\textheight}{0pt}

\renewcommand{\cftpartfont}{\bfseries}
\renewcommand{\cftpartpagefont}{\bfseries}
\setlength{\cftbeforepartskip}{1em}

\DefineVerbatimEnvironment{lstlisting}{Verbatim}{fontsize=\small,frame=single,framesep=3mm}

\newcommand{\manualnote}[1]{%
  \par\smallskip
  \noindent\begingroup
  \setlength{\fboxsep}{8pt}%
  \colorbox{undfGray}{%
    \begin{minipage}{0.95\linewidth}
      \textbf{Nota.} #1
    \end{minipage}}%
  \endgroup\par\smallskip
}

\newcommand{\manualdecision}[1]{%
  \par\smallskip
  \noindent\begingroup
  \setlength{\fboxsep}{8pt}%
  \colorbox{undfGray}{%
    \begin{minipage}{0.95\linewidth}
      \textbf{Configuração do projeto.} #1
    \end{minipage}}%
  \endgroup\par\smallskip
}

\newcommand{\manualwarning}[1]{%
  \par\smallskip
  \noindent\begingroup
  \setlength{\fboxsep}{8pt}%
  \fcolorbox{undfGreen}{white}{%
    \begin{minipage}{0.93\linewidth}
      \textbf{Atenção.} #1
    \end{minipage}}%
  \endgroup\par\smallskip
}

\newcounter{capitulosparteum}

\begin{document}
\pagenumbering{gobble}

\begin{titlepage}
  \centering
  \vspace*{1.2cm}
  {\large Universidade do Distrito Federal Professor Jorge Amaury Maia Nunes -- UnDF\par}
  \vspace{2.7cm}
  {\bfseries\Large MANUAL DE TRABALHOS DE CONCLUSÃO DE CURSO\par}
  \vspace{0.35cm}
  {\bfseries\Large E GUIA DE USO DO UnDF-TeX\par}
  \vspace{0.9cm}
  {\normalsize Orientações acadêmicas, metodológicas, normativas e técnicas\par}
  \vfill
  {\normalsize Versão do manual: 1.0.0\par}
  {\normalsize Compatível com o UnDF-TeX: 1.0.0\par}
  {\normalsize Situação: minuta para validação institucional\par}
  \medskip
  {\small Código e modelos: \url{https://github.com/igormrqs/undftex-publico}\par}
  \medskip
  {\normalsize Brasília\par}
  {\normalsize 20 de outubro de 2026\par}
\end{titlepage}

\manualfrontmatter

\chapter*{Apresentação}

Este volume reúne, em uma publicação única, o \emph{Manual de Trabalhos de Conclusão de Curso da UnDF} e o \emph{Guia de Uso do UnDF-TeX}. A Parte I orienta o percurso acadêmico, metodológico e normativo, independentemente da ferramenta de escrita. A Parte II apresenta a preparação, a edição, a compilação e a verificação técnica do template.

As duas partes permanecem disponíveis como documentos autônomos. Nesta edição unificada, porém, a capa, o sumário, a paginação e o controle documental são comuns a todo o volume.

O repositório público de código e modelos é
\url{https://github.com/igormrqs/undftex-publico}. Ele reúne a base estável de
monografia e sua documentação. Os perfis experimentais permanecem em
desenvolvimento privado; as orientações desta edição podem anteceder sua
disponibilização pública. Nesta edição, o manual e o UnDF-TeX estão
identificados pela versão 1.0.0. A cópia pública foi preparada a partir de uma
base anterior; confira seu \texttt{VERSION.txt} antes de adotar o pacote.

\caixaorientacao{A paginação editorial desta publicação não constitui modelo de paginação para o TCC. Para a contagem e a exibição das páginas do trabalho acadêmico, consulte a seção específica sobre paginação e o respectivo exemplo.}

\caixaatencao{Esta versão é uma minuta para validação institucional. Regulamentos de curso, deliberações de colegiado, normas de comitês de ética e atos oficiais prevalecem quando estabelecerem requisitos específicos.}

\clearpage
\tableofcontents
\clearpage

\manualmainmatter

\part{Manual de Trabalhos de Conclusão de Curso da \marcaUnDF}

\phantomsection
\addcontentsline{toc}{chapter}{Apresentação da Parte I}
\chapter*{Apresentação}

Este manual orienta a elaboração de Trabalhos de Conclusão de Curso no âmbito da Universidade do Distrito Federal Professor Jorge Amaury Maia Nunes -- UnDF. Sua finalidade é apoiar estudantes, docentes, orientadores, membros de banca e coordenações de curso na construção de trabalhos acadêmicos claros, rigorosos, éticos e formalmente padronizados.

O documento combina orientação acadêmica, critérios gerais de escrita científica, organização metodológica e normalização. Ele não é um manual de software e não tem como foco ensinar ferramentas de edição. Quando menciona padrões visuais, elementos obrigatórios ou decisões de apresentação, a intenção é registrar o padrão institucional que deve orientar a versão final do trabalho.

A normalização adotada considera as normas ABNT vigentes utilizadas pelo projeto e as decisões institucionais validadas durante a construção do modelo de TCC da UnDF. Nos pontos em que a ABNT admite escolhas, o manual apresenta o padrão definido e incorporado aos artefatos do projeto.

\caixaatencao{Este manual não substitui regulamentos específicos de curso, deliberações de colegiado, normas de comitês de ética ou exigências formais publicadas pela instituição. Quando houver regra específica do curso, ela deve ser observada em conjunto com este manual.}

\section*{Objetivos do manual}

São objetivos deste manual:

\begin{itemize}
  \item apresentar o sentido acadêmico do TCC e seu papel na formação superior;
  \item orientar a escolha do tema, a definição do problema e a construção dos objetivos;
  \item explicar a estrutura esperada de trabalhos acadêmicos;
  \item diferenciar elementos obrigatórios, opcionais e condicionais;
  \item orientar a escrita acadêmica, a integridade autoral e o uso responsável de fontes;
  \item esclarecer critérios gerais de apresentação, normalização, entrega e avaliação;
  \item oferecer roteiros, checklists e instrumentos de apoio ao estudante.
\end{itemize}

\section*{Como ler este manual}

O estudante que está iniciando o TCC deve ler os capítulos sobre princípios, percurso, tema, problema e projeto. O estudante em fase de escrita deve priorizar os capítulos sobre estrutura, elementos textuais, escrita acadêmica, citações e referências. O estudante em fase de entrega deve consultar os capítulos sobre normalização, avaliação, integridade, entrega e checklists.

O orientador pode usar este manual como referência comum para alinhar expectativas com seus orientandos. A banca pode utilizá-lo como apoio para verificar se a forma do trabalho está coerente com o padrão institucional, sem confundir questões de forma com a avaliação do mérito acadêmico.

\chapter{Contexto Institucional e Alcance do Manual}
\label{chap:contexto}

A Universidade do Distrito Federal Professor Jorge Amaury Maia Nunes -- UnDF é uma fundação pública integrante da administração indireta do Distrito Federal, com autonomia acadêmica, administrativa e financeira nos termos de sua legislação. A Universidade oferece formação superior pública e gratuita e articula ensino, pesquisa, extensão, inovação e relação com a comunidade em diferentes áreas do conhecimento.

A relação atualizada de cursos, unidades acadêmicas, modalidades e regulamentos específicos deve ser consultada no portal e nos atos oficiais da UnDF. O manual não fixa número de cursos nem substitui informações acadêmicas sujeitas a atualização.

\section{Base institucional}

Este manual deve ser interpretado em conjunto com a Lei Complementar distrital nº 987/2021, o Decreto distrital nº 42.333/2021, o Estatuto, o Regimento Geral, o Regimento da Graduação aprovado pela Instrução nº 24/2025, a Resolução nº 12/2023 sobre elaboração e reestruturação de PPCs e os atos posteriores aplicáveis.

\section{Regras gerais e regras de curso}

Este documento apresenta orientações gerais. Cada PPC ou regulamento pode definir modalidades admitidas, carga horária, pré-requisitos, condições de orientação, composição da banca, critérios de avaliação, forma de defesa, prazos, entrega e depósito. Em caso de conflito, deve-se consultar a coordenação e a instância institucional competente.

\regracurso

\section{Independência da ferramenta de edição}

O presente documento explica o percurso do TCC, a estrutura do trabalho, a escrita, a metodologia e a normalização sem exigir conhecimento de programas de edição.

As orientações acadêmicas permanecem válidas independentemente da ferramenta autorizada para redigir e diagramar o trabalho. Recursos automáticos podem auxiliar a aplicação do padrão, mas não alteram os requisitos acadêmicos, a responsabilidade autoral nem as regras do curso.

\fluxopendente

Os materiais editáveis de apoio devem ser publicados em canal institucional controlado, enquanto este manual permanece centrado nas orientações acadêmicas. A localização, a unidade responsável e o fluxo de manutenção ainda devem ser confirmados.

\section{Como reconhecer o tipo de orientação}

Ao longo do manual, as orientações são identificadas conforme sua origem: requisito normativo; recomendação normativa; decisão institucional validada no âmbito do projeto; regra específica de curso ou PPC; boa prática acadêmica ou editorial; e fluxo institucional ainda pendente de definição. Essa distinção evita atribuir à ABNT escolhas da UnDF ou de cada curso e impede que decisões já consolidadas sejam tratadas novamente como propostas.

\chapter{Princípios do Trabalho de Conclusão de Curso}

O Trabalho de Conclusão de Curso é uma atividade acadêmica integradora. Ele demonstra que o estudante consegue mobilizar conhecimentos construídos ao longo da formação, formular uma questão relevante, utilizar método adequado, dialogar com fontes qualificadas e comunicar resultados de modo claro e verificável.

O TCC não deve ser entendido apenas como requisito burocrático para conclusão do curso. Ele é uma oportunidade de consolidar autonomia intelectual, amadurecer práticas de investigação, desenvolver escrita acadêmica e produzir contribuição pertinente para a área de formação.

\section{Finalidade acadêmica}

A finalidade acadêmica do TCC é evidenciar competências de análise, síntese, argumentação e comunicação científica. Um bom trabalho não se limita a reunir informações: ele organiza um percurso de investigação e mostra como as evidências, leituras e escolhas metodológicas sustentam uma resposta ao problema proposto.

Entre as competências esperadas, destacam-se:

\begin{itemize}
  \item delimitar tema e problema de forma viável;
  \item localizar, selecionar e interpretar fontes confiáveis;
  \item construir objetivos coerentes com o problema;
  \item justificar escolhas metodológicas;
  \item apresentar resultados de modo organizado;
  \item discutir achados à luz da literatura;
  \item reconhecer limites da pesquisa;
  \item redigir com clareza, precisão e responsabilidade autoral.
\end{itemize}

\section{Relação entre forma e conteúdo}

A forma acadêmica existe para favorecer a leitura, a avaliação e a circulação do conhecimento. Padronização, sumário, citações, referências, paginação e elementos pré-textuais não são enfeites: eles tornam o trabalho identificável, verificável e recuperável.

No entanto, a forma não substitui o conteúdo. Um trabalho visualmente correto pode ser frágil se não apresentar problema claro, método adequado e discussão consistente. Do mesmo modo, um trabalho com boa pesquisa pode perder força se estiver desorganizado, sem referências corretas ou com apresentação descuidada.

\caixaorientacao{Durante a escrita, revise sempre duas camadas: a camada acadêmica, que envolve argumento, método e evidências; e a camada formal, que envolve estrutura, normalização e apresentação.}

\section{Autonomia e orientação}

O TCC é responsabilidade do estudante, mas não é um trabalho solitário. A orientação docente é parte essencial do processo. O orientador acompanha decisões de tema, recorte, método, bibliografia, cronograma, escrita e revisão, mas não substitui a autoria do estudante.

Cabe ao estudante comparecer às orientações preparado, registrar encaminhamentos, cumprir prazos, apresentar versões parciais e solicitar ajuda quando houver dificuldade real. Cabe ao orientador oferecer direção acadêmica, indicar problemas, sugerir caminhos e avaliar a maturidade do trabalho.

\section{Qualidade esperada}

Um TCC de qualidade apresenta coerência entre problema, objetivos, método, análise e conclusão. Também mostra domínio básico da literatura da área, escrita cuidadosa, referências adequadas e respeito às normas institucionais.

Qualidade não significa complexidade excessiva. Muitos bons trabalhos respondem a perguntas delimitadas, com método simples e bem executado. A ambição do trabalho deve ser compatível com o tempo, a formação do estudante, o acesso a dados e as exigências do curso.


\chapter{Percurso do TCC}

O TCC se desenvolve em etapas. Embora cada curso possa organizar o processo de modo próprio, há um percurso comum: escolha do tema, definição do problema, elaboração do projeto, desenvolvimento da pesquisa, redação, revisão, submissão à banca e entrega final.

\section{Etapas recomendadas}

\begin{longtable}{p{0.23\linewidth}p{0.34\linewidth}p{0.32\linewidth}}
\toprule
Etapa & Produto esperado & Pergunta orientadora \\
\midrule
\endfirsthead
\toprule
Etapa & Produto esperado & Pergunta orientadora \\
\midrule
\endhead
Exploração inicial & Lista de interesses e temas possíveis & O que vale a pena estudar na minha área? \\
Delimitação & Tema recortado e justificativa inicial & O tema é viável no tempo disponível? \\
Problematização & Problema de pesquisa & Que pergunta o trabalho pretende responder? \\
Planejamento & Projeto de pesquisa ou plano de trabalho & Como a pergunta será respondida? \\
Execução & Leituras, coleta, análise ou desenvolvimento & O método está sendo seguido com rigor? \\
Redação & Seções textuais completas e revisadas & O texto comunica o percurso de modo claro? \\
Normalização & Versão conforme padrão institucional & A forma ajuda a leitura e a avaliação? \\
Defesa ou avaliação & Apresentação e arguição, quando houver & O estudante domina o próprio trabalho? \\
Versão final & Arquivo final corrigido & As correções foram incorporadas? \\
\bottomrule
\end{longtable}

\section{Cronograma}

O cronograma deve ser realista. Trabalhos acadêmicos costumam atrasar quando o estudante subestima leitura, coleta de dados, aprovação ética, análise, revisão textual e normalização. É recomendável reservar tempo para imprevistos e para versões sucessivas.

Um cronograma mínimo deve prever:

\begin{itemize}
  \item definição do tema e aceite de orientação;
  \item elaboração ou revisão do projeto;
  \item levantamento bibliográfico;
  \item procedimentos éticos, quando aplicáveis;
  \item coleta, produção ou seleção de materiais;
  \item análise;
  \item redação das seções textuais;
  \item revisão de citações e referências;
  \item revisão formal e normalização;
  \item envio à banca;
  \item ajustes posteriores à avaliação.
\end{itemize}

\caixaatencao{O prazo de entrega à banca não deve ser tratado como prazo de término da escrita. A versão enviada para avaliação precisa chegar madura, revisada e normalizada.}

\section{Registro do processo}

O estudante deve manter registro organizado das decisões tomadas durante o TCC. Isso inclui versões do texto, referências consultadas, critérios de busca, bases utilizadas, instrumentos de coleta, pareceres, autorizações e orientações recebidas.

Esse registro facilita a escrita da metodologia, reduz risco de perda de informação e ajuda a responder perguntas da banca. Também permite reconstruir o percurso caso o trabalho precise ser revisado meses depois.

\section{Relação com a banca}

A banca avalia o trabalho entregue, não a intenção do estudante. Por isso, escolhas importantes precisam estar explícitas no texto. Se o método teve limitação, ela deve ser reconhecida. Se o recorte foi reduzido por viabilidade, isso deve ser justificado. Se determinada fonte foi privilegiada, o critério deve ser apresentado.

Uma boa defesa, quando prevista pelo curso, não compensa uma versão escrita frágil; mas uma apresentação clara pode demonstrar domínio, maturidade e capacidade de argumentação.

\chapter{Modalidades de TCC}

Os cursos podem admitir diferentes modalidades de Trabalho de Conclusão de Curso, conforme o projeto pedagógico e o regulamento específico. A presença de uma modalidade neste manual não constitui autorização automática: a decisão final cabe ao curso.

\regracurso

\section{Modalidades admitidas}

As modalidades a seguir são referências para planejamento. Em todas elas, o trabalho deve demonstrar problema ou demanda claramente definidos, fundamentação, método, análise e contribuição acadêmica compatíveis com a formação do curso.

\subsection{Monografia}

A monografia é adequada quando a investigação exige desenvolvimento extenso da fundamentação, do método e da análise. Pode ser bibliográfica, documental, empírica, aplicada ou mista.

Sua estrutura costuma reunir elementos pré-textuais, introdução, referencial teórico ou revisão de literatura, metodologia, resultados e discussão, parte conclusiva, referências e, quando necessários, apêndices e anexos. Os títulos e a divisão do desenvolvimento devem acompanhar o objeto e a tradição da área.

\subsection{Artigo científico}

O artigo científico exige recorte preciso, unidade e redação concisa. Pode ser adotado quando o curso pretende aproximar o estudante da comunicação em periódicos ou eventos, sem reduzir o rigor esperado do TCC.

Sua organização costuma incluir título, autoria, resumo, palavras-chave, introdução, método, resultados, discussão, conclusão e referências. Resultados e discussão podem aparecer juntos ou separados, conforme a área e as regras do curso ou do veículo escolhido.

Quando o artigo constituir o formato de entrega, devem ser observadas a ABNT NBR 6022:2018 e as regras do curso. A eventual adoção das instruções editoriais de um periódico não dispensa os elementos institucionais exigidos para depósito, salvo autorização expressa.

\subsection{Relatório técnico ou científico}

O relatório é adequado para intervenção, diagnóstico, prática profissional, avaliação institucional, desenvolvimento tecnológico, processo ou serviço. Ele não deve se limitar a narrar atividades.

O texto deve apresentar contexto, problema, objetivos, procedimentos, resultados, análise, recomendações quando cabíveis, parte conclusiva e referências. A descrição do que foi feito precisa ser acompanhada da justificativa das escolhas e da avaliação do que foi alcançado.

Quando aplicável, a apresentação do relatório deve considerar a ABNT NBR 10719:2015, além das exigências do PPC ou regulamento.

\subsection{Produto técnico, tecnológico ou educacional}

Essa modalidade pode envolver protocolo, aplicativo, material didático, guia, cartilha, instrumento, oficina, proposta de intervenção ou outra solução aplicada. O produto deve ser acompanhado de texto acadêmico que explique o problema, o público, os fundamentos, os requisitos, o processo de desenvolvimento, a validação e os limites de aplicação.

\subsection{Portfólio acadêmico ou de projetos}

Quando previsto pelo curso, o portfólio pode reunir produções articuladas que evidenciem percurso formativo, decisões, resultados e reflexão crítica. Ele não é uma simples coleção de arquivos: deve possuir critérios de seleção, organização, contextualização e análise das evidências apresentadas.

\subsection{Revisão de literatura}

A revisão pode constituir modalidade de TCC quando o objetivo é mapear, sintetizar ou analisar criticamente a produção existente. Deve explicitar bases consultadas, estratégia de busca, período, critérios de inclusão e exclusão e procedimento de análise. Uma sequência de resumos sem método e sem síntese crítica não caracteriza revisão acadêmica suficiente.

\subsection{Estudo de caso}

O estudo de caso pode ser admitido como formato de trabalho ou empregado como estratégia de pesquisa em outra modalidade. Deve delimitar o caso, justificar sua escolha, identificar as fontes de evidência e desenvolver análise que ultrapasse o relato descritivo.

\section{Escolha da modalidade}

A modalidade deve ser escolhida com o orientador e confirmada conforme o regulamento do curso. A decisão deve considerar a pergunta de pesquisa, os objetivos, a natureza dos resultados, o tempo disponível, o acesso a dados, as exigências éticas e a forma de avaliação. A modalidade organiza a entrega acadêmica; ela não substitui o método da pesquisa.

\caixaatencao{Pôster e slides são formas de apresentação do trabalho, não modalidades do documento de TCC.}

As orientações para pôsteres e slides devem ser publicadas como material de apoio à apresentação. Caso a UnDF adote pôster científico institucional, a norma específica vigente deve ser conferida antes da publicação do modelo.

\chapter{Tema, Problema, Objetivos e Justificativa}

A escolha do tema é uma das decisões mais importantes do TCC. Um tema adequado precisa ser relevante, viável, delimitado e compatível com a área de formação. A partir dele, o estudante constrói problema, objetivos e justificativa.

\section{Escolha do tema}

O tema pode surgir de experiências profissionais, disciplinas cursadas, estágios, projetos de extensão, iniciação científica, observação de problemas sociais, lacunas na literatura ou demandas institucionais. O interesse pessoal é importante, mas não basta. O tema precisa permitir investigação acadêmica.

Antes de decidir, verifique:

\begin{itemize}
  \item se há bibliografia suficiente e atual;
  \item se o recorte cabe no tempo disponível;
  \item se há acesso aos dados, documentos, campo ou participantes;
  \item se o orientador possui aderência ao tema;
  \item se há necessidade de aprovação ética;
  \item se o tema dialoga com a formação do curso.
\end{itemize}

\section{Delimitação}

Delimitar é transformar um assunto amplo em objeto investigável. Um tema como ``saúde mental de estudantes'' é amplo demais. Um recorte mais viável poderia considerar um curso, uma instituição, um período, um grupo específico, uma intervenção, uma política ou uma dimensão do fenômeno.

A delimitação pode envolver:

\begin{itemize}
  \item espaço: instituição, território, serviço, comunidade ou ambiente;
  \item tempo: período histórico, ano, semestre, ciclo ou marco específico;
  \item população: grupo, perfil, categoria profissional ou público-alvo;
  \item material: documentos, bases, obras, sistemas, processos ou práticas;
  \item enfoque: dimensão teórica, metodológica, política, técnica ou social.
\end{itemize}

\section{Problema de pesquisa}

O problema de pesquisa é a pergunta central do trabalho. Ele deve ser claro, investigável e coerente com o método. Um problema muito amplo leva a objetivos genéricos. Um problema muito restrito pode impedir análise significativa.

Perguntas de pesquisa costumam usar verbos e expressões como: de que modo, em que medida, quais fatores, quais efeitos, como ocorre, quais relações, que percepções, que desafios, que contribuições.

\caixaorientacao{Um bom teste é perguntar: ao final do trabalho, será possível responder a essa pergunta com base no método escolhido? Se a resposta for não, o problema ainda precisa ser ajustado.}

\section{Objetivo geral e objetivos específicos}

O objetivo geral expressa a finalidade principal do trabalho. Os objetivos específicos indicam etapas necessárias para alcançar o objetivo geral. Eles não devem ser uma lista de tarefas mecânicas, mas ações intelectuais ligadas à investigação.

Verbos úteis incluem analisar, compreender, descrever, identificar, comparar, avaliar, investigar, discutir, mapear, propor, sistematizar e interpretar. Verbos vagos, como conhecer ou falar sobre, devem ser evitados.

\section{Justificativa}

A justificativa explica por que o trabalho deve ser realizado. Ela pode apresentar relevância acadêmica, social, institucional, profissional, científica ou tecnológica. Também pode indicar lacunas, demandas, problemas práticos ou necessidade de sistematização.

Uma justificativa consistente não depende de afirmações grandiosas. Ela demonstra, com argumentos e fontes, que o tema merece atenção e que o recorte proposto é pertinente.


\chapter{Projeto de Pesquisa ou Plano de Trabalho}

O projeto é o documento de planejamento do TCC. Ele organiza as decisões antes da execução e permite que orientador, estudante e curso avaliem viabilidade, coerência e riscos.

Sua apresentação deve considerar a ABNT NBR 15287:2025 e as exigências do curso. O projeto não é uma versão abreviada do TCC final: ele descreve uma pesquisa planejada, enquanto o trabalho final registra e analisa o percurso efetivamente realizado.

\section{Função do projeto}

O projeto serve para responder antecipadamente a quatro perguntas: o que será estudado, por que será estudado, como será estudado e em quanto tempo será realizado. Ele não é uma formalidade isolada; é a base do trabalho final.

Mesmo quando o curso não exige entrega formal de projeto, o estudante deve elaborar um plano mínimo. Trabalhos iniciados sem planejamento tendem a enfrentar problemas de recorte, método e prazo.

O projeto não é uma versão abreviada do TCC final. Ele registra o que se pretende fazer e pode ser ajustado de forma justificada durante a execução. O trabalho final, por sua vez, relata o percurso efetivamente realizado, analisa os resultados e apresenta a parte conclusiva.

\section{Organização do projeto}

Conforme a ABNT NBR 15287:2025, o projeto é organizado em elementos pré-textuais, textuais e pós-textuais. O conteúdo e a extensão variam conforme a natureza da pesquisa e as regras do curso.

\subsection{Elementos pré-textuais}

A folha de rosto é obrigatória no projeto. Capa, listas e sumário são usados conforme a necessidade e a orientação institucional. O modelo do curso deve indicar quais informações de identificação e vínculo acadêmico serão exigidas.

\subsection{Elementos textuais}

Um projeto de TCC pode conter:

\begin{itemize}
  \item título provisório;
  \item tema e delimitação;
  \item problema de pesquisa;
  \item hipótese ou pressuposto, quando couber;
  \item objetivo geral e objetivos específicos;
  \item justificativa;
  \item referencial teórico inicial;
  \item metodologia;
  \item cronograma;
  \item cronograma.
\end{itemize}

As seções não precisam receber exatamente esses títulos, mas o texto deve permitir reconhecer tema, problema, objetivos, justificativa, base teórica, método, recursos quando necessários e planejamento temporal.

\subsection{Elementos pós-textuais}

As referências são obrigatórias. Glossário, apêndices, anexos e índice são incluídos quando necessários. Instrumentos de coleta elaborados pelo estudante, como roteiros e questionários, são apêndices; documentos de terceiros são anexos.

\section{Referencial teórico inicial}

O referencial teórico inicial apresenta autores, conceitos e debates que sustentam a pesquisa. Não precisa esgotar a literatura, mas deve demonstrar que o estudante conhece o campo mínimo necessário para iniciar o trabalho.

É recomendável evitar duas fragilidades: listar autores sem articulação ou transformar o referencial em resumo de obras. O referencial deve construir base conceitual para o problema.

\section{Metodologia no projeto}

A metodologia deve explicar o tipo de pesquisa, os materiais ou participantes, os procedimentos, os critérios de análise e os cuidados éticos. Quanto mais claro for o método no projeto, menor a chance de retrabalho durante a execução.

Quando houver coleta com pessoas, dados sensíveis, intervenção, prontuários, imagens, entrevistas ou observação de campo, a necessidade de avaliação ética deve ser verificada com antecedência.

\section{Cronograma e riscos}

O cronograma deve considerar etapas reais e riscos previsíveis. Se o trabalho depende de autorização externa, resposta de participantes, acesso a documentos ou aprovação ética, essas etapas precisam aparecer no planejamento.

\caixaatencao{Não deixe revisão bibliográfica, normalização e conferência de referências para a última semana. Essas etapas costumam revelar inconsistências que exigem reescrita.}

\section{Recursos e viabilidade}

Quando necessários, devem ser descritos recursos humanos, materiais, tecnológicos e financeiros, bem como formas de acesso a campo, bases, equipamentos ou documentos. Dependências críticas precisam aparecer no planejamento e no cronograma.

\chapter{Ética na Pesquisa e Responsabilidade Acadêmica}

A ética na pesquisa envolve respeito às pessoas, aos animais, às fontes, aos dados, às instituições e à autoria. Ela deve orientar o desenho do estudo desde o projeto, e não ser acrescentada apenas como declaração formal na versão final.

\section{Pesquisas com seres humanos}

Entrevistas, questionários, grupos focais, observação, intervenções e acesso a prontuários, imagens, gravações ou dados pessoais podem caracterizar pesquisa envolvendo seres humanos. O enquadramento deve ser verificado com o orientador e com a instância institucional competente antes do recrutamento ou da coleta.

Quando a apreciação pelo Sistema CEP/Conep for aplicável, o protocolo deve ser apresentado com a documentação exigida e a pesquisa somente pode começar após a aprovação ética. Nos TCCs de graduação abrangidos pela Resolução CNS nº 510/2016, o registro no Comitê de Ética em Pesquisa ocorre sob responsabilidade do orientador. Pesquisa vinculada a projeto do orientador já aprovado pode exigir emenda, conforme o caso.

Nem toda atividade acadêmica com pessoas possui o mesmo enquadramento. Atividades exclusivamente educacionais, levantamentos sem finalidade de pesquisa e uso de informações de acesso público, por exemplo, devem ser analisados de acordo com a regulamentação aplicável; o estudante não deve presumir dispensa por conta própria.

\caixaatencao{Consulte o orientador e a instância competente antes de recrutar participantes, acessar prontuários ou iniciar qualquer coleta. O manual não autoriza dispensa de apreciação ética.}

\subsection{Documentos e autorizações}

Conforme o desenho da pesquisa, podem ser necessários projeto detalhado, instrumentos de coleta, termos de consentimento ou assentimento, autorização do local ou da instituição, autorização para acesso a registros e parecer ético. Esses documentos têm funções diferentes: a autorização institucional não substitui a aprovação ética, e nenhuma delas substitui o processo de consentimento dos participantes quando ele for aplicável.

\section{Pesquisas com animais e material biológico}

Atividades de ensino ou pesquisa científica com animais devem observar a Lei nº 11.794/2008, as normas do Conselho Nacional de Controle de Experimentação Animal e os procedimentos da Comissão de Ética no Uso de Animais. Quando exigida, a autorização da CEUA deve ser obtida antes do início da atividade.

Pesquisas com material biológico humano, dados clínicos ou procedimentos de saúde podem exigir cuidados adicionais. O estudante deve verificar o fluxo institucional e a regulamentação específica da área ainda na fase de projeto.

\section{Consentimento, assentimento e privacidade}

O consentimento livre e esclarecido é um processo de comunicação. O participante deve compreender objetivos, procedimentos, riscos, benefícios, voluntariedade, formas de registro e uso das informações. Quando houver crianças, adolescentes ou pessoas sem plena capacidade de consentir, deve-se verificar também a necessidade de assentimento e de autorização do responsável.

A forma de registro do consentimento depende do contexto e das normas aplicáveis. Assinatura em papel não é a única possibilidade em todos os estudos, mas qualquer alternativa deve ser justificada e aprovada quando a pesquisa estiver sujeita à apreciação ética.

\section{Proteção, anonimização e gestão dos dados}

O projeto deve definir quais dados são realmente necessários, quem terá acesso, onde serão armazenados, por quanto tempo serão mantidos e como serão eliminados ou compartilhados. Dados pessoais e sensíveis exigem medidas proporcionais ao risco, conforme a legislação e as orientações institucionais.

O plano de gestão deve registrar finalidade, base legal a ser confirmada institucionalmente, responsáveis, controle de acesso, cópias de segurança, retenção, descarte e condições de compartilhamento. Dados não devem ser enviados a serviços externos ou ferramentas de inteligência artificial sem autorização e avaliação do risco de exposição.

Anonimização e pseudonimização não são sinônimos. Na pseudonimização, os registros ainda podem ser associados a uma pessoa mediante informação adicional mantida separadamente. A anonimização busca impedir a reidentificação por meios razoáveis, mas deve ser tratada como processo baseado em risco, e não como garantia absoluta.

Imagens, voz, gravações e combinações de atributos podem permitir reidentificação mesmo após a retirada do nome. A divulgação deve considerar esse risco, o consentimento obtido, a finalidade informada e eventuais limites éticos ou jurídicos.

\subsection{Separação entre validação e respostas}

Em coletas digitais, pode ser necessário confirmar elegibilidade ou impedir respostas duplicadas sem associar a identidade ao conteúdo respondido. Uma arquitetura possível é validar o participante em uma etapa separada, emitir um código de uso único e armazenar as respostas sem identificadores diretos. A tabela de correspondência, se indispensável, deve ficar separada, protegida e acessível apenas a responsáveis autorizados.

\section{Validade e autenticidade da coleta}

A proteção da identidade deve conviver com a qualidade da pesquisa. O projeto precisa definir critérios de inclusão, prevenção de duplicidade, teste do instrumento, registro de recusas e perdas e procedimentos para identificar respostas inválidas sem expor participantes.

Se a pesquisa pretende fazer inferências sobre uma população, a estratégia amostral e seus limites devem ser explícitos. Quantidade de respostas não corrige uma amostra enviesada. Controles de autenticidade, quando necessários, devem ser proporcionais ao risco e informados aos participantes; vigilância excessiva pode violar privacidade e comprometer a própria pesquisa.

\section{Ética em pesquisa bibliográfica e documental}

Pesquisas bibliográficas e documentais também exigem ética. O estudante deve representar corretamente as ideias dos autores, evitar distorções, citar as fontes consultadas e distinguir análise própria de informação obtida em documentos.

Fontes digitais devem ser avaliadas com cuidado. Sempre que possível, devem ser priorizados livros, artigos científicos, legislação, documentos institucionais, bases oficiais e publicações técnicas reconhecidas.

\section{Plágio, citação e crédito acadêmico}

Plágio é apresentar texto, ideia, dado, imagem, estrutura ou resultado alheio como se fosse próprio. A cópia literal sem indicação, a paráfrase sem referência e a tradução não identificada omitem a autoria e violam a integridade acadêmica.

Citar não é plagiar. O estudante pode reproduzir um trecho nos limites pertinentes ou explicar uma ideia com suas próprias palavras, desde que identifique a fonte de acordo com as regras de citação e inclua a referência correspondente. A citação dá crédito; ela não autoriza montagem excessiva de texto alheio nem substitui a análise do autor do TCC.

\caixaorientacao{Ao usar um trabalho anterior, deixe visíveis três camadas: o que pertence à fonte, onde a fonte pode ser localizada e qual é a interpretação ou contribuição própria do TCC. O capítulo sobre integridade acadêmica detalha as formas de plágio e sua prevenção.}

Para os tipos de plágio, o uso responsável de tecnologias e as medidas de prevenção, consulte o Capítulo \ref{chap:integridade}.

\section{Conflitos de interesse}

Quando o estudante pesquisa seu próprio local de trabalho, estágio, serviço ou instituição, deve reconhecer possíveis conflitos de interesse e vieses. A proximidade com o objeto pode enriquecer a pesquisa, mas exige transparência, proteção contra coerção e cuidado na interpretação.

\caixaorientacao{A ética precisa aparecer no desenho da pesquisa, na coleta, na análise, na escrita, no armazenamento e na divulgação dos resultados.}

\section{Retenção, descarte e compartilhamento}

O protocolo e o termo de consentimento, quando aplicáveis, devem ser coerentes sobre prazo de guarda, possibilidade de reutilização e compartilhamento. O descarte precisa ser seguro e verificável. Antes de depositar bases junto ao TCC, confirme se a divulgação é autorizada e se os dados estão adequadamente anonimizados.

\section{Três autorizações com funções diferentes}

A autorização da instituição pesquisada permite acesso ao local, aos registros ou à atividade, dentro de seus limites. O consentimento ou assentimento protege a decisão livre do participante. A aprovação ética avalia o protocolo no sistema competente. Um documento não substitui automaticamente os demais.

\chapter{Metodologia de Pesquisa}

A metodologia explica o caminho utilizado para produzir conhecimento. Ela deve ser compatível com o problema, os objetivos e a modalidade do TCC. Um método bem descrito permite que o leitor compreenda como as informações foram obtidas, organizadas e analisadas.

\section{Abordagens de pesquisa}

A abordagem qualitativa busca compreender sentidos, experiências, processos, relações, discursos e contextos. A abordagem quantitativa trabalha com mensuração, frequência, associação, comparação estatística ou descrição numérica. A abordagem mista combina procedimentos qualitativos e quantitativos.

A escolha da abordagem não deve ser feita por preferência pessoal, mas pela natureza do problema. Perguntas sobre percepções, significados e processos costumam exigir abordagem qualitativa. Perguntas sobre prevalência, frequência e associação podem exigir abordagem quantitativa. Problemas complexos podem se beneficiar de desenho misto.

\section{Pesquisa bibliográfica}

A pesquisa bibliográfica utiliza produção acadêmica já publicada como base principal. Deve indicar bases, descritores, critérios de seleção, período, tipos de documento e estratégia de análise.

Uma pesquisa bibliográfica frágil apenas reúne textos encontrados. Uma pesquisa bibliográfica consistente seleciona fontes de modo justificado, compara argumentos e constrói síntese crítica.

\section{Pesquisa documental}

A pesquisa documental analisa documentos que podem incluir legislação, relatórios, atas, políticas públicas, registros institucionais, processos, planos, diretrizes, imagens, materiais de divulgação, dados administrativos ou arquivos históricos.

O estudante deve explicar origem dos documentos, critérios de seleção, confiabilidade, contexto de produção e forma de análise.

\section{Pesquisa de campo}

A pesquisa de campo envolve aproximação direta com pessoas, grupos, instituições ou ambientes. Pode utilizar entrevistas, questionários, observação, grupos focais, registros fotográficos, diário de campo ou instrumentos específicos da área.

Esse tipo de pesquisa exige planejamento cuidadoso, autorização quando necessária, atenção ética e registro organizado.

\section{População, amostra e corpus}

O trabalho deve identificar de onde provêm as evidências analisadas. Em pesquisas com pessoas, isso inclui população de interesse, critérios de inclusão e exclusão, forma de recrutamento, tamanho e composição da amostra. Em pesquisas documentais ou bibliográficas, inclui a delimitação do corpus, das bases, do período e dos tipos de documento.

Quando houver intenção de generalizar resultados, a estratégia amostral precisa sustentar essa pretensão. Amostras por conveniência podem ser úteis, mas suas limitações devem ser reconhecidas. Distribuição etária, territorial ou de outros grupos relevantes deve ser comparada ao público sobre o qual se pretende concluir.

\section{Entrevistas}

Entrevistas podem ser estruturadas, semiestruturadas ou abertas. O roteiro deve estar alinhado aos objetivos e evitar perguntas que induzam respostas.

Na análise, o estudante não deve apenas transcrever falas. É necessário interpretar os dados, organizar categorias, relacionar respostas ao problema e dialogar com a literatura.

\section{Questionários}

Questionários podem coletar dados de modo padronizado. Devem usar perguntas claras, linguagem adequada ao público, alternativas coerentes e ordem lógica. Antes da aplicação, recomenda-se testar o instrumento com pequeno grupo.

Dados de questionários precisam ser apresentados e analisados com cuidado. Percentuais sem contexto podem gerar interpretações equivocadas, especialmente em amostras pequenas.

\section{Validade e controle de qualidade da coleta}

O planejamento deve prever teste do instrumento, treinamento da equipe quando necessário, registro de recusas e perdas, prevenção de duplicidade e critérios para tratar respostas incompletas, inconsistentes ou suspeitas. Em formulários digitais, controles contra automação e múltiplos envios podem ser necessários, desde que não coletem identificadores excessivos.

Autenticidade não significa necessariamente conhecer a identidade civil do participante. O objetivo é reunir evidências suficientes de elegibilidade e unicidade de acordo com o risco do estudo, preservando a privacidade. Procedimentos de validação devem ser descritos na metodologia e compatíveis com o protocolo ético.

\section{Observação}

A observação pode ser participante ou não participante, estruturada ou aberta. Deve indicar local, período, foco de observação, forma de registro e cuidados éticos.

O diário de campo é recurso importante para registrar situações, impressões, eventos e reflexões. Esses registros devem ser tratados como material de análise, não como memória informal.

\section{Análise de dados}

A análise transforma material bruto em resultados interpretáveis. Pode envolver categorização, análise temática, análise de conteúdo, estatística descritiva, comparação documental, análise de discurso, triangulação ou outro procedimento adequado.

O estudante deve explicar como passou dos dados aos resultados. A banca precisa enxergar o caminho analítico.

\section{Plano de gestão dos dados}

Pesquisas empíricas devem indicar como os dados serão nomeados, armazenados, protegidos, versionados e preparados para análise. Quando houver dados pessoais, o plano também deve definir perfis de acesso, separação de identificadores, prazo de retenção e forma de descarte ou compartilhamento.

O conjunto bruto não deve ser alterado silenciosamente. Exclusões, correções, transformações e critérios de tratamento precisam ser registrados para que o percurso analítico possa ser auditado.

\section{Limitações metodológicas}

Todo trabalho possui limites. Reconhecê-los não enfraquece o TCC; ao contrário, demonstra maturidade. Limitações podem envolver amostra, acesso a dados, tempo, escopo, instrumentos, generalização ou disponibilidade de fontes.

\caixaatencao{Método não é decoração textual. Se a metodologia não explica como o trabalho foi realizado, o leitor não consegue avaliar a validade dos resultados.}

\chapter{Estrutura Geral do TCC}

O TCC deve ser organizado em elementos pré-textuais, textuais e pós-textuais. Essa organização facilita a identificação do trabalho, a leitura do conteúdo, a localização de partes específicas e a conferência das fontes utilizadas.

\section{Grupos de elementos}

\begin{longtable}{p{0.24\linewidth}p{0.34\linewidth}p{0.30\linewidth}}
\toprule
Grupo & Função & Exemplos \\
\midrule
\endfirsthead
\toprule
Grupo & Função & Exemplos \\
\midrule
\endhead
Pré-textuais & Identificar o trabalho e orientar a consulta inicial & Capa, folha de rosto, resumo, listas e sumário \\
Textuais & Desenvolver a investigação e apresentar a argumentação acadêmica & Parte introdutória, desenvolvimento e parte conclusiva \\
Pós-textuais & Registrar fontes e complementar o conteúdo & Referências, glossário, apêndices e anexos \\
\bottomrule
\end{longtable}

\section{Elementos obrigatórios e opcionais}

Alguns elementos são obrigatórios na versão final. Outros são opcionais e devem ser usados quando contribuírem para a identificação, a consulta ou a compreensão do trabalho.

\requisitonormativo

\begin{longtable}{p{0.38\linewidth}p{0.22\linewidth}p{0.28\linewidth}}
\toprule
Elemento & Situação geral & Observação \\
\midrule
\endfirsthead
\toprule
Elemento & Situação geral & Observação \\
\midrule
\endhead
Capa & Obrigatório & Identifica instituição, autoria, título, local e ano. \\
Folha de rosto & Obrigatório & Traz natureza do trabalho, curso e orientação. \\
Dados internacionais de catalogação & Condicionado ao fluxo institucional & A norma prevê sua posição, mas a emissão e a exigência devem seguir a orientação da Biblioteca Central da UnDF. \\
Errata & Opcional & Usada apenas quando há correções posteriores. \\
Folha de aprovação & Obrigatório & Registra aprovação e composição avaliadora. \\
Dedicatória & Opcional & Texto breve e pessoal. \\
Agradecimentos & Opcional & Reconhece contribuições relevantes. \\
Epígrafe & Opcional & Citação breve relacionada ao trabalho. \\
Resumo em português & Obrigatório & Síntese objetiva do trabalho. \\
Resumo em língua estrangeira & Obrigatório & Usualmente em inglês, mas o idioma pode variar por decisão do curso. \\
Listas & Opcional & Usadas quando há quantidade relevante de itens. \\
Sumário & Obrigatório & Último elemento pré-textual. \\
Introdução & Obrigatório & Apresenta tema, problema, objetivos e percurso. \\
Desenvolvimento & Obrigatório & Contém revisão, método, análise e discussão. \\
Parte conclusiva & Obrigatório & Pode receber o título ``Conclusão'' ou ``Considerações finais'', conforme a área e a decisão do curso. \\
Referências & Obrigatório & Lista as fontes citadas. \\
Glossário & Opcional & Define termos técnicos. \\
Apêndice & Opcional & Material elaborado pelo autor. \\
Anexo & Opcional & Material não elaborado pelo autor. \\
Índice & Opcional & Organiza termos de forma sistemática e segue a ABNT NBR 6034:2004. \\
\bottomrule
\end{longtable}

\section{Ordem dos elementos}

A ordem dos elementos deve ser estável. Alterar a posição de resumo, listas, sumário, referências, apêndices ou anexos compromete a padronização e pode dificultar a avaliação.

O sumário deve aparecer ao final dos elementos pré-textuais. A parte textual começa depois do sumário. Referências e demais pós-textuais aparecem após as considerações finais.

A ABNT NBR 14724:2024 define a macroestrutura textual, mas deixa a nomenclatura dos títulos a critério do autor. O trabalho deve ser organizado em seções. Assim, referencial teórico, metodologia, dados, resultados e discussão são formas possíveis de desenvolver o conteúdo, e não títulos universais impostos pela norma.

\section{Coerência estrutural}

A estrutura deve refletir o percurso do trabalho. Se o TCC é pesquisa empírica, a metodologia precisa explicar coleta e análise. Se é pesquisa bibliográfica, a metodologia precisa explicar critérios de busca e seleção. Se é desenvolvimento de produto, o texto deve apresentar requisitos, processo de construção, avaliação e limitações.

\caixaatencao{Não existe estrutura única para todos os cursos e objetos, mas existe uma exigência comum: o leitor precisa entender o caminho seguido e a relação entre problema, método, resultados e conclusão.}

\chapter{Elementos Pré-textuais}

Os elementos pré-textuais antecedem o conteúdo principal. Eles identificam o trabalho, registram informações institucionais e facilitam a navegação pelo documento. A capa não é contada. A contagem começa na folha de rosto, mas o número só aparece na primeira página textual. A página reservada aos dados internacionais de catalogação não participa da contagem.

\requisitonormativo

\section{Capa}

A capa identifica a instituição, o autor, o título, o local e o ano. O título deve representar com precisão o conteúdo do trabalho. Títulos excessivamente longos dificultam a leitura; títulos genéricos reduzem a identidade acadêmica do estudo.

No padrão institucional da UnDF, a capa deve valorizar a identidade da universidade e manter apresentação limpa. O curso não precisa aparecer na capa quando já constar na folha de rosto e na natureza do trabalho.

\section{Folha de rosto}

A folha de rosto apresenta dados essenciais para identificação acadêmica: autor, título, subtítulo quando houver, natureza do trabalho, curso, instituição, grau pretendido, orientação, local e ano.

A natureza do trabalho deve indicar que se trata de Trabalho de Conclusão de Curso apresentado ao curso correspondente, na Universidade do Distrito Federal Professor Jorge Amaury Maia Nunes -- UnDF, como requisito parcial para obtenção do grau definido.

\section{Dados internacionais de catalogação}

Os dados internacionais de catalogação devem seguir o código de catalogação vigente e o fluxo definido pela Biblioteca Central da UnDF. A ABNT NBR 14724:2024 determina sua posição no verso da folha de rosto ou, em documento eletrônico, na página imediatamente posterior, mas a exigência, a emissão e o momento de solicitação dependem da orientação institucional.

O estudante não deve criar sozinho uma ficha catalográfica nem inventar classificação, notação de autor ou descritores. A página que contém esses dados não é contada nem numerada.

\fluxopendente

Devem ser confirmados pela Biblioteca Central da UnDF o serviço responsável, os critérios de emissão, o momento da solicitação, o formato aceito e a relação com o depósito no repositório.

\section{Errata}

A errata é elemento eventual. Ela registra correções identificadas depois da finalização ou impressão do trabalho. Deve conter a referência do trabalho e a indicação objetiva de folha, linha, trecho incorreto e forma correta.

Não se deve usar errata para corrigir problemas extensos de redação, estrutura ou normalização que poderiam ter sido resolvidos antes da entrega.

\section{Folha de aprovação}

A folha de aprovação registra a avaliação do trabalho, a data e os membros da banca. A ABNT exige nome, titulação acadêmica e instituição dos integrantes. No padrão institucional validado no âmbito do projeto, a titulação aparece de forma sintética antes do nome, por exemplo, ``Prof. Dr. Nome'', sem uma segunda linha de título acadêmico. A data de aprovação segue a forma: Brasília, DF, dia de mês de ano.

Os nomes da banca devem ser apresentados de forma uniforme, com prefixo acadêmico antes do nome. A titulação não deve ser repetida de maneira extensa em linha separada.

\decisaoinstitucional

\section{Dedicatória, agradecimentos e epígrafe}

Dedicatória, agradecimentos e epígrafe são opcionais. A dedicatória é breve e pessoal. Os agradecimentos podem reconhecer orientação, apoio institucional, colaboração técnica, financiamento, familiares e pessoas que contribuíram para o percurso.

A epígrafe pré-textual é uma citação relacionada ao tema geral do trabalho, apresentada conforme o modelo gráfico institucional. Deve indicar autoria e possuir referência correspondente. A epígrafe também pode aparecer na abertura de uma seção primária; nesse caso, relaciona-se especificamente ao conteúdo da seção e continua sujeita às regras de citação e referência.

Esses elementos devem manter sobriedade acadêmica. Textos muito longos, excessivamente informais ou desconectados do contexto institucional devem ser evitados.

\section{Resumo}

O resumo apresenta uma síntese do trabalho em um único parágrafo. Deve indicar objetivo, método, resultados principais e conclusão ou considerações finais, em linguagem clara e concisa. Para trabalhos acadêmicos, a ABNT NBR 6028:2021 recomenda de 150 a 500 palavras. Ele não deve conter citações, notas, figuras, tabelas, enumerações extensas ou informações que não estejam desenvolvidas no corpo do trabalho.

As palavras-chave devem representar conceitos centrais do trabalho. Devem aparecer logo abaixo do resumo, após a expressão ``Palavras-chave:'', separadas por ponto e vírgula e encerradas por ponto. A ABNT não estabelece uma quantidade universal; use quantidade adequada ao conteúdo e escolha termos específicos que favoreçam a recuperação da informação.

\decisaoinstitucional

\section{Resumo em língua estrangeira}

O resumo em língua estrangeira corresponde ao resumo em português, vertido para outro idioma. O idioma mais comum é o inglês, mas o curso ou a instituição pode admitir outro idioma conforme a área.

O resumo estrangeiro deve preservar equivalência de conteúdo, não ser uma versão reduzida, ampliada ou divergente do resumo em português.

\caixaatencao{Resumo em português e resumo em língua estrangeira são elementos obrigatórios distintos na ABNT NBR 14724:2024. A equivalência de conteúdo não autoriza a supressão de um deles.}

\section{Listas}

Listas de ilustrações, tabelas, abreviaturas, siglas e símbolos são usadas quando ajudam o leitor a localizar itens recorrentes. Não é necessário criar lista para quantidade mínima de elementos quando isso não favorece a consulta.

\section{Sumário}

O sumário apresenta a estrutura interna do trabalho e deve refletir fielmente os títulos e páginas do texto. Elementos pré-textuais não entram no sumário. O sumário é o último elemento pré-textual.

Em arquivos digitais, recomenda-se que as entradas sejam links navegáveis. O sumário deve ser atualizado somente depois da paginação final e conferido no PDF entregue.

\section{Exemplo de contagem e exibição}

O Apêndice \ref{ap:paginacao} apresenta uma sequência completa. Em síntese: a capa fica fora da contagem; a folha de rosto inicia a contagem; a página dos dados catalográficos é desconsiderada; os demais pré-textuais continuam a contagem sem mostrar o número; e a primeira página da introdução exibe o número correspondente à sua posição.

\chapter{Elementos Textuais}

De acordo com a ABNT NBR 14724:2024, o texto é composto por uma parte introdutória, pelo desenvolvimento e por uma parte conclusiva. A norma deixa a nomenclatura dos títulos a critério do autor e orienta a organização do trabalho em seções. Portanto, a sequência apresentada neste capítulo é um modelo acadêmico recomendável, não uma lista de títulos obrigatórios para todas as áreas.

\section{Introdução}

A introdução deve situar o leitor. Ela apresenta tema, delimitação, problema, objetivos, justificativa, síntese do método e organização do trabalho.

Uma introdução eficiente responde, de modo articulado, às seguintes perguntas:

\begin{itemize}
  \item Qual é o tema e qual recorte foi adotado?
  \item Qual problema será enfrentado?
  \item Por que o problema é relevante?
  \item O que o trabalho pretende alcançar?
  \item Como a investigação foi conduzida?
  \item Como o texto está organizado?
\end{itemize}

A introdução abre o percurso e prepara o leitor para as seções seguintes. Ela não deve antecipar toda a discussão nem se transformar em revisão bibliográfica completa.

\section{Desenvolvimento}

O desenvolvimento é a parte central do trabalho. Sua divisão depende da natureza da pesquisa, da modalidade do TCC e das práticas da área. As seções a seguir representam uma organização frequente e podem ser adaptadas com o orientador.

\subsection{Referencial teórico ou revisão de literatura}

O referencial apresenta conceitos, autores e debates que sustentam a análise. Deve ser seletivo e articulado ao problema. Não basta resumir textos em sequência; é necessário construir diálogo entre as fontes, explicitar divergências e mostrar como o trabalho se posiciona.

\subsection{Metodologia}

A metodologia descreve como o trabalho foi realizado e deve permitir que o leitor avalie a adequação do percurso. Conforme o tipo de pesquisa, pode indicar:

\begin{itemize}
  \item abordagem e desenho da pesquisa;
  \item população, amostra, corpus, documentos, bases ou materiais;
  \item critérios de inclusão e exclusão;
  \item instrumentos e procedimentos de coleta;
  \item procedimentos de tratamento e análise;
  \item cuidados éticos e proteção de dados;
  \item limitações metodológicas.
\end{itemize}

\caixaorientacao{A metodologia deve ser escrita no nível de detalhe suficiente para que outro leitor compreenda o caminho percorrido. Frases genéricas como ``foi feita pesquisa na internet'' não atendem à exigência acadêmica.}

\subsection{Dados ou corpus: seção opcional}

Uma seção primária própria para dados ou corpus pode ser inserida depois da metodologia quando o conjunto analisado exigir caracterização extensa. Ela pode descrever origem, variáveis, composição da amostra, critérios de qualidade, dados ausentes, transformações, organização do corpus e produto final do tratamento.

Essa seção é opcional. Se a caracterização for breve, deve permanecer como subseção da metodologia ou ser incorporada aos resultados. Separá-la só se justifica quando melhora a compreensão e evita que a metodologia fique desproporcional. Ela não deve repetir o procedimento de coleta nem antecipar a interpretação dos achados.

\subsection{Resultados e discussão}

Os resultados apresentam o que foi encontrado, produzido ou sistematizado. A discussão interpreta esses achados à luz do problema, dos objetivos e da literatura, explicando significados, convergências, diferenças, implicações e limites.

Há duas organizações igualmente possíveis:

\begin{itemize}
  \item \textbf{Resultados e discussão integrados:} cada achado é apresentado e imediatamente interpretado. Essa forma costuma favorecer estudos qualitativos, aplicados ou organizados por temas.
  \item \textbf{Resultados e discussão separados:} primeiro são apresentados os achados e depois é feita a interpretação conjunta. Essa forma pode ser útil quando há grande volume de dados ou tradição disciplinar específica.
\end{itemize}

A escolha deve ser feita com o orientador. Resultados não devem ser acumulados sem explicação, e a discussão não deve introduzir dados que não foram apresentados.

\section{Conclusão ou considerações finais}

A parte conclusiva retoma o problema e os objetivos, sintetiza os principais achados e explicita a resposta que o trabalho alcançou. Também pode registrar contribuições, aplicações, limitações e recomendações ou estudos futuros coerentes com os resultados.

Os títulos ``Conclusão'' e ``Considerações finais'' são admissíveis quando compatíveis com a orientação do curso. Em geral, ``conclusão'' comunica um fechamento mais direto, enquanto ``considerações finais'' pode ser apropriado quando o objeto admite síntese reflexiva sem pretensão de encerramento definitivo.

Essa parte não deve apresentar teoria nova, dados inéditos ou discussão que não tenha sido desenvolvida anteriormente. Um bom fechamento responde ao que o trabalho prometeu investigar, sem repetir mecanicamente o resumo ou a introdução.

\chapter{Elementos Pós-textuais}

Os elementos pós-textuais completam o trabalho após a parte textual. Eles registram fontes, esclarecem termos, apresentam materiais complementares e preservam documentos necessários à compreensão da pesquisa.

\section{Referências}

As referências listam as fontes citadas no texto. Devem seguir a norma vigente adotada pela instituição e manter consistência de autoria, título, edição, local, editora, ano, link e data de acesso quando aplicável.

No modelo de TCC do UnDF-TeX, inclusive no formato de artigo, as referências mantêm a apresentação do TCC convencional: abertura em nova página, título centralizado sem indicativo numérico, entradas alinhadas à margem esquerda, espaçamento simples dentro de cada referência e separação uniforme entre elas. A camada bibliográfica é compartilhada pelos perfis. Essa padronização da entrega acadêmica não altera quais elementos são obrigatórios ou opcionais; uma futura submissão a revista seguirá suas próprias instruções editoriais.

O estudante deve conferir se toda fonte citada aparece nas referências e se toda referência listada foi efetivamente citada no texto, quando a orientação do curso exigir lista apenas de fontes citadas.

\section{Glossário}

O glossário reúne termos técnicos ou especializados acompanhados de definição. Ele é útil quando o trabalho utiliza vocabulário específico que pode dificultar a leitura por público não especializado.

Não deve ser usado para definir termos comuns ou para substituir explicações conceituais que deveriam aparecer no referencial teórico.

\section{Apêndices}

Apêndices são materiais elaborados pelo autor do trabalho. Podem incluir instrumentos de coleta, roteiros de entrevista, questionários, quadros analíticos, produtos educacionais, materiais de apoio, documentos produzidos na pesquisa ou dados complementares.

Apêndices devem ser citados no texto. Não se deve incluir material sem relação direta com o trabalho apenas para aumentar volume.

\section{Anexos}

Anexos são materiais não elaborados pelo autor, mas necessários à compreensão, comprovação ou contextualização do trabalho. Podem incluir leis, documentos institucionais, mapas, formulários externos, autorizações, protocolos ou imagens de fonte externa.

Assim como os apêndices, anexos devem ser citados no texto e ter relação clara com a pesquisa.

\section{Índice}

O índice é elemento opcional e pouco comum em TCCs. Quando utilizado, deve permitir localização de termos, nomes ou assuntos relevantes. Em trabalhos de graduação, normalmente o sumário e as listas são suficientes.

\caixaatencao{Material complementar não deve esconder fragilidades do texto. Se uma informação é essencial para entender o argumento, ela precisa ser discutida no corpo do trabalho, ainda que o documento completo apareça em apêndice ou anexo.}


\chapter{Escrita Acadêmica}

A escrita acadêmica deve ser clara, objetiva, precisa e responsável. Seu objetivo é comunicar uma investigação, não impressionar pelo excesso de termos difíceis. Um texto acadêmico de qualidade permite que o leitor compreenda o problema, acompanhe o método, avalie os argumentos e verifique as fontes.

\section{Clareza e precisão}

Clareza não significa simplificação excessiva. Significa escolher palavras adequadas, organizar ideias em ordem lógica e evitar ambiguidades. Precisão significa usar conceitos de modo consistente, definir termos importantes e evitar generalizações sem evidência.

Evite expressões vagas, como ``muitos autores dizem'', ``a sociedade precisa'', ``desde sempre'', ``é evidente'' ou ``todos sabem''. Afirmações acadêmicas precisam de delimitação, fonte ou argumento.

\section{Coesão e progressão}

O texto deve avançar. Cada parágrafo precisa contribuir para a construção do argumento. Um parágrafo acadêmico geralmente apresenta uma ideia central, desenvolve essa ideia e se conecta ao parágrafo seguinte.

Problemas comuns de progressão incluem:

\begin{itemize}
  \item parágrafos que apenas acumulam citações;
  \item repetição da mesma ideia com palavras diferentes;
  \item mudança brusca de assunto;
  \item conceitos usados antes de serem definidos;
  \item resultados apresentados antes do método;
  \item conclusão que não responde aos objetivos.
\end{itemize}

\section{Voz autoral}

Ter voz autoral não significa escrever em tom pessoal ou opinativo. Significa conduzir o texto, selecionar fontes, estabelecer relações, interpretar dados e assumir responsabilidade pelas conclusões.

O estudante deve evitar dois extremos: desaparecer atrás de citações ou escrever sem diálogo com a literatura. Um bom TCC mostra que o autor leu, compreendeu e usou fontes para construir análise própria.

\section{Uso de primeira pessoa}

A aceitação de primeira pessoa varia por área e por orientação institucional. Em muitos trabalhos, recomenda-se linguagem impessoal ou primeira pessoa do plural de modo controlado. O mais importante é manter consistência e sobriedade.

Independentemente da escolha, o texto deve evitar informalidade, subjetivismo excessivo e frases que fragilizem a análise, como ``eu acho'', ``na minha opinião'' ou ``acredito que'', quando não houver base argumentativa.

\section{Revisão textual}

Revisar não é apenas corrigir ortografia. A revisão deve observar coerência, ordem das seções, repetição, citações, referências, tabelas, figuras, títulos, paginação e conformidade com o padrão institucional.

Uma estratégia eficiente é revisar em camadas:

\begin{enumerate}
  \item coerência entre problema, objetivos, método e conclusão;
  \item estrutura das seções e progressão argumentativa;
  \item parágrafos, frases e transições;
  \item citações, referências, notas e fontes;
  \item normalização, sumário, listas e paginação;
  \item leitura final do PDF.
\end{enumerate}

\caixaorientacao{Sempre que possível, deixe intervalo entre escrita e revisão. O estudante identifica mais problemas quando relê o texto depois de alguma distância.}

\chapter{Normalização e Apresentação Gráfica}

A normalização assegura unidade visual, facilita a leitura e permite que os trabalhos sejam avaliados por critérios comuns. No padrão da UnDF, as decisões de apresentação consideram as normas ABNT aplicáveis e escolhas institucionais para os pontos em que há margem de definição.

\section{Papel, margens e espaçamento}

O trabalho deve ser apresentado em papel A4. O padrão institucional utiliza margem superior e esquerda de 3 cm, margem inferior e direita de 2 cm. O corpo textual utiliza espaçamento de uma linha e meia, com exceções para elementos que exigem espaçamento simples, como citações longas, notas, referências, natureza do trabalho e fontes de ilustrações.

\requisitonormativo

O formato A4 e as margens indicadas seguem os requisitos de apresentação da ABNT NBR 14724:2024. O espaçamento de uma linha e meia no corpo acompanha a recomendação normativa e foi adotado pelo projeto. O recuo de parágrafo e a ausência de espaço adicional são decisões institucionais validadas no âmbito do projeto.

\recomendacaonormativa

\decisaoinstitucional

\section{Fonte e legibilidade}

A ABNT recomenda tamanho 12 para o texto e tamanho menor e uniforme para exceções, mas não determina uma família tipográfica. Para compatibilizar os modelos editáveis, o projeto adota o seguinte padrão institucional:

\decisaoinstitucional

\begin{longtable}{p{0.42\linewidth}p{0.21\linewidth}p{0.25\linewidth}}
\toprule
Elemento & Tamanho & Tratamento \\
\midrule
Corpo do texto & 12 pt & Família serifada institucional \\
Seção primária & 12 pt & Negrito e caixa alta \\
Demais seções & 12 pt & Negrito e diferenciação hierárquica consistente \\
Citação direta longa & 10 pt & Recuo de 4 cm, espaço simples, sem aspas \\
Notas de rodapé & 10 pt & Espaço simples \\
Legendas e fontes & 10 pt & Espaço simples \\
Natureza do trabalho & 10 pt & Espaço simples \\
Conteúdo de quadros e tabelas & 10 pt ou 12 pt & Uniforme e legível no documento \\
\bottomrule
\end{longtable}

\fluxopendente

A família tipográfica institucional ainda precisa de deliberação. Até que exista padrão aprovado, o estudante deve seguir a orientação vigente do curso e manter uma única família legível e uniforme em todo o trabalho, sem combinar fontes apenas por efeito visual.

Mais importante do que escolher uma fonte incomum é manter consistência. Misturar muitas fontes, tamanhos e estilos prejudica a aparência institucional do trabalho.

\section{Parágrafos}

O padrão institucional adota recuo de 1,25 cm nos parágrafos, inclusive no primeiro parágrafo após títulos. O texto principal não deve usar espaçamento extra entre parágrafos como substituto do recuo, salvo em elementos específicos definidos pelo modelo.

\section{Paginação}

A capa não apresenta número visível. A folha de rosto é contada, mas a numeração visível aparece a partir da parte textual. O número deve permanecer em posição estável e seguir a contagem até os elementos pós-textuais.

A página dos dados internacionais de catalogação não entra na contagem. No anverso, o número aparece no canto superior direito, a 2 cm das bordas superior e direita. O Apêndice \ref{ap:paginacao} mostra a sequência completa.

\section{Títulos e numeração progressiva}

Os títulos devem refletir a hierarquia do trabalho. Seções primárias, secundárias, terciárias, quaternárias e quinárias devem ser numeradas de forma progressiva, respeitando a organização do conteúdo.

O padrão institucional atual apresenta as seções primárias no formato ``1 INTRODUÇÃO'', sem a palavra ``Capítulo'', com número e título na mesma linha e título em caixa alta.

\decisaoinstitucional

O uso de caixa alta na seção primária e os destaques gráficos dos demais níveis são escolhas da UnDF. Pela ABNT NBR 6024:2012, a numeração deve limitar-se até a seção quinária e não se usa ponto, hífen, travessão ou outro sinal entre o indicativo e o título.

Títulos sem indicativo numérico, como resumo, listas, sumário, referências, glossário, apêndices, anexos e índice, recebem tratamento próprio. Elementos pré-textuais não entram no sumário; referências e demais elementos pós-textuais entram sem indicativo numérico.

\section{Alíneas e subalíneas}

Alíneas são usadas quando o texto precisa apresentar enumeração interna sem transformar cada item em seção. Devem manter paralelismo gramatical: se o primeiro item começa com verbo, os demais também devem seguir estrutura semelhante; se começa com substantivo, os demais devem manter o padrão.

Use alíneas com moderação. Listas muito longas podem indicar que o conteúdo deveria ser reescrito em parágrafos, quadro ou tabela.

\section{Siglas}

Na primeira ocorrência, a sigla deve aparecer após a forma por extenso. Depois disso, a sigla pode ser usada sozinha.

Exemplo: Universidade do Distrito Federal Professor Jorge Amaury Maia Nunes (UnDF). Nas ocorrências seguintes, pode-se usar UnDF.

Siglas pouco conhecidas, siglas numerosas ou siglas específicas de área podem justificar uma lista de abreviaturas e siglas.

\section{Abreviaturas e meses}

Abreviaturas devem seguir uso padronizado e não devem prejudicar a clareza. Em referências, datas e meses precisam respeitar o padrão normativo adotado. O estudante deve evitar misturar meses por extenso, abreviados e em outro idioma na mesma lista de referências.

\section{Impressão e versão digital}

Mesmo quando a entrega é digital, a apresentação deve considerar leitura em tela e eventual impressão. Figuras, tabelas, gráficos e cores devem permanecer legíveis. Evite depender exclusivamente de cor para transmitir informação, pois a impressão em preto e branco pode prejudicar a compreensão.

\section{Sumário}

O sumário deve reproduzir os títulos e a paginação do trabalho. Não deve incluir elementos pré-textuais. Deve ser conferido depois da versão final, pois alterações de página e título podem desatualizar suas entradas.

\section{Consistência}

A normalização exige consistência. Se uma figura recebe fonte abaixo, todas devem receber. Se a banca usa determinado formato, todos os membros devem seguir o mesmo padrão. Se as referências usam um padrão de data de acesso, ele deve ser mantido.

\caixaatencao{A revisão da versão de entrega deve ser feita no PDF, não apenas no arquivo editável. Alguns problemas de quebra de página, paginação, figuras e sumário só aparecem claramente no documento gerado ou exportado.}

\chapter{Modelos Comentados dos Elementos}

Este capítulo apresenta modelos esquemáticos dos principais elementos do TCC. Os modelos não substituem o padrão gráfico oficial, mas ajudam o estudante a compreender a função de cada parte e a conferir se os dados essenciais foram preenchidos.

\section{Capa}

\modeloinstitucional{MODELO ESQUEMÁTICO DE CAPA}{
\centering
UNIVERSIDADE DO DISTRITO FEDERAL PROFESSOR JORGE AMAURY MAIA NUNES -- UnDF\par
\vspace{2em}
NOME DO AUTOR\par
\vspace{4em}
TÍTULO DO TRABALHO\par
Subtítulo, se houver\par
\vfill
BRASÍLIA\par
ANO\par
}

A capa deve identificar o trabalho de modo direto. O título precisa ser fiel ao conteúdo e não deve prometer mais do que a pesquisa entrega. O subtítulo, quando usado, deve complementar o título, não repeti-lo.

Erros comuns: título genérico, nome do autor diferente da folha de rosto, ano desatualizado, excesso de informação na capa e uso de siglas não explicadas.

\section{Folha de rosto}

\modeloinstitucional{MODELO ESQUEMÁTICO DE FOLHA DE ROSTO}{
\centering
NOME DO AUTOR\par
\vspace{3em}
TÍTULO DO TRABALHO\par
Subtítulo, se houver\par
\vspace{3em}
\raggedleft
\begin{minipage}{0.58\linewidth}
Trabalho de Conclusão de Curso apresentado ao curso de [CURSO] da Universidade do Distrito Federal Professor Jorge Amaury Maia Nunes -- UnDF, como requisito parcial para obtenção do título de [GRAU].\par
\vspace{0.5em}
Orientador: Nome do Orientador.
\end{minipage}\par
\vfill
\centering
BRASÍLIA\par
ANO\par
}

A folha de rosto repete dados centrais e acrescenta a natureza do trabalho. É nela que o leitor identifica o curso, o grau pretendido e a orientação.

\section{Folha de aprovação}

\modeloinstitucional{MODELO ESQUEMÁTICO DE FOLHA DE APROVAÇÃO}{
\centering
NOME DO AUTOR\par
\vspace{2em}
TÍTULO DO TRABALHO\par
\vspace{2em}
\raggedright
Trabalho aprovado como requisito parcial para obtenção do título de [GRAU] no curso de [CURSO] da Universidade do Distrito Federal Professor Jorge Amaury Maia Nunes -- UnDF.\par
\vspace{1.2em}
Brasília, DF, dia de mês de ano.\par
\vspace{2em}
\centering
\rule{0.68\linewidth}{0.4pt}\par
Prof. Dr. Nome do Orientador\par
Instituição\par
\vspace{1.5em}
\rule{0.68\linewidth}{0.4pt}\par
Profa. Dra. Nome do Membro\par
Instituição\par
}

A folha de aprovação deve manter a banca em formato uniforme. Quando o padrão institucional adota titulação abreviada antes do nome, ela não precisa ser repetida em linha separada.

\section{Errata}

\modeloinstitucional{MODELO ESQUEMÁTICO DE ERRATA}{
\centering
ERRATA\par
\vspace{1em}
\raggedright
SOBRENOME, Nome. Título do trabalho. Ano. Trabalho de Conclusão de Curso (Graduação em Curso) -- Universidade do Distrito Federal Professor Jorge Amaury Maia Nunes -- UnDF, Brasília, ano.\par
\vspace{1em}
\centering
\begin{tabular}{cccc}
\toprule
Folha & Linha & Onde se lê & Leia-se \\
\midrule
16 & 10 & auto-clavado & autoclavado \\
\bottomrule
\end{tabular}
}

A errata deve ser usada apenas para correções pontuais. Se o problema exige reescrever seções inteiras, o correto é revisar a versão final antes da entrega.

\section{Resumo}

\modeloinstitucional{MODELO ESQUEMÁTICO DE RESUMO}{
\centering
RESUMO\par
\vspace{1em}
\raggedright
Texto conciso apresentando objetivo, método, resultados e conclusões do trabalho. O resumo deve permitir que o leitor compreenda a essência da pesquisa sem consultar o texto integral.\par
\vspace{1em}
\textbf{Palavras-chave:} termo 1; termo 2; termo 3.
}

O resumo deve evitar citações, notas, ilustrações e informações não desenvolvidas no trabalho. As palavras-chave precisam representar conceitos centrais.

\section{Sumário}

\modeloinstitucional{MODELO ESQUEMÁTICO DE SUMÁRIO}{
\raggedright
1 INTRODUÇÃO \dotfill 10\par
2 REFERENCIAL TEÓRICO \dotfill 15\par
3 METODOLOGIA \dotfill 28\par
4 RESULTADOS E DISCUSSÃO \dotfill 35\par
5 CONSIDERAÇÕES FINAIS \dotfill 48\par
REFERÊNCIAS \dotfill 52\par
APÊNDICE A -- Roteiro de entrevista \dotfill 56\par
ANEXO A -- Documento institucional \dotfill 58\par
}

O sumário deve refletir exatamente os títulos e páginas da versão final. Elementos pré-textuais não entram no sumário.

\section{Referências}

\modeloinstitucional{MODELO ESQUEMÁTICO DE REFERÊNCIAS}{
\raggedright
SOBRENOME, Nome. \textbf{Título da obra}. Edição. Local: Editora, ano.\par
\vspace{0.5em}
SOBRENOME, Nome. Título do artigo. \textbf{Título do periódico}, local, volume, número, páginas, ano. Disponível em: endereço. Acesso em: dia mês ano.\par
}

As referências devem ser conferidas uma a uma. A lista final deve ser coerente com as citações do texto.

\section{Apêndice e anexo}

\modeloinstitucional{MODELO ESQUEMÁTICO DE APÊNDICE E ANEXO}{
\centering
APÊNDICE A -- ROTEIRO DE ENTREVISTA\par
\vspace{2em}
\raggedright
Material elaborado pelo autor do trabalho.\par
\vspace{2em}
\centering
ANEXO A -- DOCUMENTO INSTITUCIONAL\par
\vspace{2em}
\raggedright
Material não elaborado pelo autor, mas necessário à compreensão do trabalho.
}

Apêndices e anexos devem ser citados no corpo do texto. Material não citado tende a parecer acréscimo sem função acadêmica.


\chapter{Citações e Referências}

Citações e referências são parte central da integridade acadêmica. Elas mostram de onde vêm informações, conceitos, dados e argumentos utilizados no trabalho. Também permitem que o leitor verifique as fontes e compreenda o diálogo do TCC com a produção existente.

\section{Quando citar}

Deve-se citar sempre que o texto utilizar ideia, conceito, dado, interpretação, trecho literal, imagem, tabela, método, classificação ou argumento de outra fonte. A citação é necessária mesmo quando o estudante reescreve a ideia com suas próprias palavras.

Não é necessário citar conhecimentos amplamente consolidados e gerais, mas, em caso de dúvida, é mais seguro indicar a fonte que sustentou a afirmação.

\section{Citação indireta}

A citação indireta apresenta a ideia de outro autor com as palavras do estudante. Ela deve preservar o sentido original e indicar a fonte. Parafrasear não significa trocar algumas palavras; significa compreender a ideia e reescrevê-la de modo próprio.

\section{Citação direta curta}

A citação direta curta reproduz literalmente trecho de uma fonte e permanece incorporada ao parágrafo, entre aspas, com indicação de autoria, ano e página ou outro localizador disponível. Deve ser usada com moderação, apenas quando a formulação original for importante.

\section{Citação direta longa}

A citação direta longa, com mais de três linhas, é destacada do parágrafo, sem aspas, em espaço simples e com fonte menor. A ABNT recomenda recuo de 4 cm; a UnDF adota esse valor como decisão institucional validada no âmbito do projeto. A citação deve indicar página ou outro localizador disponível e ser comentada pelo estudante. Não se deve inserir bloco citado sem explicar sua relevância para o argumento.

\section{Citação de citação}

A citação de citação deve ser excepcional e evitada quando a fonte original puder ser consultada. Na chamada, identifica-se a autoria original e a fonte efetivamente consultada com a expressão \emph{apud}; na lista de referências, registra-se somente a obra consultada. O ideal é que o estudante consulte diretamente as obras centrais de seu tema.

\section{Localizadores e documentos não paginados}

Nas citações diretas, a página deve ser indicada quando existir. Em documentos sem paginação, use seção, capítulo, artigo, parágrafo, posição, timestamp ou outro identificador capaz de levar o leitor ao trecho. Na citação indireta, a indicação de página ou localização é opcional, mas pode ser útil em obras extensas.

\section{Sistema autor-data}

Nas chamadas entre parênteses, sobrenomes de pessoas e nomes de entidades são grafados em letras maiúsculas e minúsculas, por exemplo, (Silva, 2024) e (Associação Brasileira de Normas Técnicas, 2024). Siglas consolidadas permanecem em caixa alta, como (IBGE, 2025). A abreviação \emph{et al.} pode ser utilizada em obras com quatro ou mais autores e deve seguir padrão uniforme em todo o trabalho.

\section{Referências}

A lista de referências deve permitir identificação e localização das fontes. Obras impressas, artigos, capítulos, legislação, documentos digitais, sites, relatórios, vídeos e outros materiais exigem informações específicas.

Para situações especiais de chamada e localização, consulte o Capítulo \ref{chap:casos-citacao}. Para modelos completos de referências por tipo documental, consulte o Capítulo \ref{chap:modelos-referencias}.

Erros comuns em referências incluem:

\begin{itemize}
  \item autor citado no texto ausente da lista final;
  \item referência listada sem citação no texto;
  \item links quebrados ou sem data de acesso quando necessária;
  \item títulos incompletos;
  \item autores grafados de formas diferentes;
  \item mistura de estilos;
  \item ausência de informações essenciais.
\end{itemize}

\caixaatencao{A conferência cruzada é obrigatória: toda fonte citada deve aparecer nas referências, e toda referência deve corresponder a uma fonte efetivamente utilizada no trabalho.}

\section{Fontes confiáveis}

Nem toda fonte é adequada para TCC. O estudante deve priorizar literatura científica, livros acadêmicos, documentos oficiais, legislação, relatórios técnicos e bases reconhecidas. Fontes jornalísticas, blogs e conteúdos de redes sociais podem ser úteis em objetos específicos, mas precisam de justificativa e análise crítica.

\chapter{Casos Especiais de Citação}
\label{chap:casos-citacao}

As citações devem permitir que o leitor identifique a fonte utilizada e compreenda a relação entre a ideia citada e o argumento do trabalho. Este capítulo apresenta situações frequentes que exigem atenção.

Os nomes, datas, títulos e trechos utilizados nos exemplos deste capítulo são fictícios e têm finalidade exclusivamente didática, salvo quando uma fonte real for expressamente identificada. Autores fictícios usados apenas para demonstrar chamadas não exigem inclusão na lista final de referências.

\section{Um autor}

Quando a obra possui um autor, a chamada deve indicar sobrenome e ano, conforme o sistema adotado. A citação pode aparecer integrada ao texto ou entre parênteses.

Exemplo de chamada integrada: Segundo Silva (2024), a delimitação do problema orienta as escolhas metodológicas.

Exemplo de chamada parentética: A delimitação do problema orienta as escolhas metodológicas (Silva, 2024).

\section{Dois ou três autores}

Quando houver dois ou três autores, todos devem ser indicados na chamada, respeitando a forma adotada pela norma e pelo padrão institucional.

Exemplo: A escrita acadêmica exige planejamento, revisão e domínio das fontes (Silva; Almeida; Costa, 2024).

\section{Mais de três autores}

Em obras com mais de três autores, pode-se indicar o primeiro autor seguido da expressão adequada, conforme a norma adotada. O estudante deve manter o mesmo padrão em todo o trabalho.

Exemplo: A revisão bibliográfica deve seguir critérios explícitos de seleção (Silva et al., 2024).

\section{Autor entidade}

Quando a autoria é de instituição, órgão, associação, universidade ou entidade coletiva, a chamada deve identificar a entidade responsável. Siglas podem ser usadas quando forem conhecidas e estiverem apresentadas adequadamente.

Exemplo: A apresentação de trabalhos acadêmicos deve seguir critérios de estrutura e normalização (Associação Brasileira de Normas Técnicas, 2024).

\section{Fonte sem autoria identificada}

Quando não houver autor identificado, a chamada deve ser feita pelo título ou pela primeira palavra significativa do título, conforme a regra aplicável. Fontes sem autoria devem ser usadas com cuidado e avaliadas quanto à confiabilidade.

\section{Autores com mesmo sobrenome}

Quando autores diferentes possuem o mesmo sobrenome e publicações no mesmo ano ou no mesmo contexto de citação, é necessário acrescentar elementos que permitam diferenciá-los, como iniciais dos prenomes, conforme o sistema adotado.

\section{Mais de uma obra do mesmo autor}

Quando o mesmo autor possui mais de uma obra no mesmo ano, usam-se letras após o ano para diferenciar as obras, desde que a lista de referências também mantenha essa distinção.

Exemplo: Silva (2024a) discute a delimitação do problema, enquanto Silva (2024b) aborda a construção dos objetivos.

\section{Citação de várias obras}

Quando uma afirmação é sustentada por várias obras, as fontes devem ser organizadas conforme o padrão normativo adotado. O estudante deve evitar acúmulo desnecessário de referências e selecionar as mais pertinentes.

Exemplo: A revisão de literatura deve ser seletiva e articulada ao problema (Almeida, 2022; Costa, 2023; Silva, 2024).

\section{Citação direta com supressão}

Quando parte do trecho citado é omitida, a supressão deve ser indicada de forma padronizada. A omissão não pode alterar o sentido original.

Exemplo: ``A pesquisa acadêmica [...] exige coerência entre problema, método e análise'' (Silva, 2024, p. 25).

\section{Citação direta com acréscimo}

Quando o estudante acrescenta palavra ou explicação dentro de citação direta, deve marcar que o acréscimo não pertence ao texto original.

Exemplo: ``A instituição [universidade] deve definir critérios de apresentação'' (Silva, 2024, p. 30).

\section{Citação direta com grifo}

Quando houver destaque em trecho citado, deve-se indicar se o grifo pertence ao original ou se foi acrescentado pelo autor do TCC.

Exemplo: ``O método deve ser compatível com o problema'' (Silva, 2024, p. 41, grifo nosso).

\section{Citação traduzida}

Quando o estudante traduz trecho de obra em outro idioma, deve informar que a tradução é sua, conforme o padrão adotado.

Exemplo: ``A revisão deve ser transparente e reprodutível'' (Smith, 2023, p. 15, tradução nossa).

\section{Informação verbal}

Informações obtidas oralmente, como palestras, entrevistas informais ou comunicações pessoais, devem ser usadas com cautela. Quando forem relevantes, a origem deve ser indicada no texto ou em nota, conforme orientação institucional.

O estudante deve avaliar se a informação verbal pode ser verificada. Quando não puder, talvez seja melhor não usá-la como sustentação central do argumento.

\section{Documento não paginado}

Em documentos digitais sem paginação, a localização pode ser indicada por seção, capítulo, parágrafo, tempo de vídeo ou outra forma que ajude o leitor a localizar o trecho, quando necessário.

\section{Citação de legislação}

Na citação de legislação, deve-se identificar o ente responsável, o tipo de ato, número, data e informações necessárias para localização. No texto, evite citar apenas o número da lei sem contextualizar seu conteúdo.

\section{Citação de dados e bases}

Dados estatísticos, bases públicas e levantamentos institucionais também precisam de fonte. Informe órgão, ano, tabela, base, painel, relatório ou recorte utilizado.

\caixaorientacao{O estudante deve padronizar as citações desde o início. Corrigir dezenas de chamadas no fim do trabalho é mais trabalhoso e aumenta o risco de inconsistência.}

\chapter{Ilustrações, Tabelas e Elementos Especiais}

Ilustrações, tabelas, gráficos, quadros, fluxogramas, equações e outros elementos especiais devem servir à compreensão do trabalho. Eles não devem ser inseridos apenas para preencher páginas ou tornar o texto visualmente mais chamativo.

\section{Ilustrações}

Ilustrações incluem figuras, gráficos, mapas, fotografias, desenhos, esquemas, fluxogramas e outros recursos visuais. Devem ser citadas no texto, posicionadas próximas ao trecho em que são discutidas e acompanhadas de identificação e fonte.

A identificação deve permitir que o leitor localize e entenda o item. A fonte informa origem ou autoria da imagem. Quando a ilustração for elaborada pelo próprio estudante, isso deve ser indicado.

A identificação aparece acima do elemento, com palavra designativa, número, travessão e título, por exemplo: ``Figura 3 -- Etapas da pesquisa''. Abaixo ficam a fonte e, quando necessárias, notas explicativas. O padrão institucional para autoria própria é: ``Fonte: elaboração própria.'' Essa fórmula deve ser usada de modo uniforme em figuras, quadros, gráficos, fluxogramas e demais ilustrações produzidas pela autoria do TCC.

\decisaoinstitucional

\section{Figuras comparativas}

Quando duas imagens são usadas para comparação, cada uma deve ter relação clara com o argumento. A simples colocação lado a lado não dispensa explicação no texto. O estudante deve indicar o que deve ser observado, quais diferenças importam e como elas sustentam a análise.

Se as imagens forem figuras independentes, cada uma recebe número, título, fonte e chamada próprios, mesmo que sejam posicionadas lado a lado. Use partes ``a'' e ``b'' sob uma única legenda apenas quando elas constituírem, de fato, uma figura composta inseparável.

\section{Tabelas e quadros}

Tabelas organizam dados, especialmente numéricos. Quadros costumam organizar informações textuais, categorias, comparações ou sínteses. Ambos devem ser legíveis, citados no texto e acompanhados de fonte.

As tabelas devem observar as Normas de apresentação tabular do IBGE. Seu número e título ficam acima; fonte e notas ficam abaixo. Quadros são tratados como ilustrações no padrão do trabalho e seguem a identificação institucional correspondente.

Evite tabelas muito largas, com excesso de linhas ou informação que poderia ser apresentada em texto. Quando necessário, dados extensos podem ser deslocados para apêndice.

\subsection{Diferença prática entre tabela e quadro}

Use tabela quando o foco estiver em dados numéricos, medidas, frequências, percentuais ou resultados organizados estatisticamente. Use quadro quando o foco estiver em texto, conceitos, categorias, critérios, sínteses ou comparação qualitativa.

Essa distinção ajuda a manter o padrão de apresentação e evita que todo elemento organizado em linhas e colunas seja chamado genericamente de tabela.

\subsection{Cuidados com tabelas}

Tabelas devem apresentar cabeçalho claro, unidades de medida, fonte e alinhamento que favoreça comparação. Não repita no texto todos os números da tabela; interprete os resultados mais importantes.

Quando uma tabela continuar em outra página, repita o cabeçalho e indique a continuidade com as expressões ``continua'', ``continuação'' e ``conclusão'', conforme a posição. Evite quebrar linhas de dados entre páginas. Unidades de medida devem estar no cabeçalho ou em posição que elimine ambiguidade.

\subsection{Cuidados com quadros}

Quadros devem ser usados para síntese textual. Evite quadros com parágrafos muito longos ou conteúdo que seria melhor explicado no corpo do texto.

\section{Gráficos}

Gráficos ajudam a visualizar tendências, proporções e comparações. Devem ter título claro, eixos identificados, unidades quando aplicáveis e fonte. Cores e estilos devem ser escolhidos para favorecer legibilidade, inclusive em impressão em preto e branco.

Não use apenas cor para distinguir séries. Combine cor com padrões, marcadores, rótulos ou estilos de linha e preserve contraste suficiente. Quando o arquivo for digital, forneça texto alternativo ou descrição equivalente no documento acessível.

\section{Fluxogramas}

Fluxogramas representam processos, etapas, decisões ou fluxos de trabalho. Devem usar símbolos de modo consistente e indicar início, fim, atividades e decisões quando representarem um processo completo.

No padrão institucional, fluxogramas podem ser tratados como ilustrações quando aparecem em pequena quantidade. Quando forem recorrentes no trabalho, pode ser adotada lista própria para facilitar consulta.

A ABNT NBR 14724:2024 disciplina a apresentação do fluxograma como ilustração, mas não fixa uma gramática completa de símbolos. O trabalho deve adotar convenção reconhecida na área, explicá-la quando necessário e manter consistência. Símbolos, setas e decisões precisam ser legíveis e não podem depender exclusivamente de cores.

\section{Equações}

Equações devem ser numeradas quando forem citadas posteriormente ou quando houver necessidade de referência. Variáveis e símbolos precisam ser definidos no texto ou em lista própria, se o trabalho utilizar muitos símbolos.

\section{Pseudocódigo}

Pseudocódigo pode ser usado em trabalhos que descrevem algoritmos, rotinas, processos computacionais ou procedimentos lógicos. Deve ser legível para público acadêmico e não depender de detalhes de uma linguagem de programação específica, salvo quando esse for o objeto do trabalho.

Apresente o pseudocódigo como algoritmo ou elemento técnico numerado, com título, fonte, entradas, saídas e passos claramente identificados. Numeração de linhas e palavras-chave em destaque ajudam a discussão. Se o bloco não couber na página, prefira movê-lo integralmente, reduzir espaços ou dividi-lo em partes logicamente nomeadas; não reduza a fonte a ponto de prejudicar a leitura.

\caixaorientacao{Todo elemento visual deve ser comentado no texto. A frase ``conforme figura a seguir'' é insuficiente se o estudante não explicar o que a figura demonstra.}

\chapter{Orientação, Apresentação e Avaliação}

A avaliação do TCC considera mérito acadêmico, coerência metodológica, qualidade da escrita, normalização e domínio do estudante sobre o próprio trabalho. Cada curso pode definir critérios e formas de apresentação específicos.

\section{Papel do orientador}

O orientador acompanha o desenvolvimento do TCC, mas não é coautor automático do texto. Sua função é ajudar o estudante a qualificar decisões acadêmicas, corrigir rumos, indicar bibliografia, avaliar viabilidade e preparar o trabalho para avaliação.

O estudante deve respeitar os prazos combinados e enviar versões com antecedência. Orientações produtivas dependem de material concreto para leitura.

\section{Papel da banca}

A banca examina o trabalho, identifica contribuições, aponta fragilidades e avalia se o estudante atingiu os objetivos esperados. Entre os aspectos avaliáveis estão:

\begin{itemize}
  \item relevância e delimitação do tema;
  \item clareza do problema e dos objetivos;
  \item adequação metodológica e ética;
  \item qualidade da fundamentação e da análise;
  \item coerência da parte conclusiva;
  \item uso adequado de citações e referências;
  \item apresentação formal e domínio do estudante.
\end{itemize}

\section{Apresentação do trabalho}

A apresentação comunica o TCC já produzido; ela não constitui uma modalidade do documento. O curso ou o evento pode exigir defesa oral, slides, pôster ou combinação dessas formas. Tempo, dimensões, arquivo, identidade visual e critérios de avaliação devem seguir a convocação específica.

\subsection{Defesa oral e slides}

Na defesa oral, o estudante deve expor tema, problema, objetivos, método, resultados, conclusão e contribuições dentro do tempo definido. Os slides devem apoiar a fala com títulos informativos, texto enxuto e elementos visuais legíveis. A apresentação não deve tentar reproduzir todas as páginas do TCC.

É recomendável ensaiar com o arquivo final, testar fontes, vídeos e conexões, manter uma cópia em formato PDF e conferir previamente os recursos disponíveis na sala.

\subsection{Pôster técnico ou científico}

O pôster é uma forma de comunicação visual sintética, comum em seminários e congressos. Deve permitir leitura rápida e apoiar a explicação presencial do estudante. Em geral, apresenta:

\begin{itemize}
  \item título, autoria, orientação e vínculo institucional;
  \item contexto ou introdução breve;
  \item objetivo;
  \item método;
  \item resultados principais;
  \item conclusão;
  \item referências essenciais e informações de contato, quando solicitadas.
\end{itemize}

Gráficos, tabelas, fluxos e imagens devem ter função informativa, fonte identificada e resolução suficiente. Excesso de texto, fonte pequena e baixo contraste prejudicam a leitura. O modelo e as dimensões do evento prevalecem sobre recomendações genéricas.

\section{Arguição}

Durante a arguição, o estudante deve ouvir atentamente, anotar recomendações e responder com clareza. Quando não souber responder, é melhor reconhecer a limitação do que improvisar justificativas frágeis.

Discordâncias acadêmicas podem ocorrer. O estudante pode defender suas escolhas com respeito, base argumentativa e abertura à revisão.

\section{Correções após a banca}

Quando a banca indicar correções, o estudante deve incorporá-las na versão final dentro do prazo estabelecido. A versão completa deve ser revisada novamente, pois alterações pontuais podem afetar sumário, paginação, referências e coerência textual.

\chapter{Acessibilidade dos Trabalhos Acadêmicos}
\label{chap:acessibilidade}

A acessibilidade permite que o trabalho seja consultado por pessoas com diferentes condições visuais, auditivas, motoras, cognitivas ou tecnológicas. Enquanto a UnDF não publicar requisitos próprios para TCC acessível, as orientações deste capítulo constituem boas práticas editoriais e devem ser aplicadas sempre que a ferramenta permitir.

\boapratica

\section{Documento digital}

O PDF deve conter texto pesquisável, títulos marcados em hierarquia lógica, ordem de leitura coerente, idioma definido, links descritivos e propriedades básicas do documento. Sempre que possível, gere PDF marcado, com estrutura reconhecível por tecnologias assistivas, em vez de digitalizar páginas como imagens.

Use os estilos de título da ferramenta de edição. A aparência visual isolada, como aumentar e colocar o texto em negrito, não cria necessariamente a estrutura que leitores de tela usam para navegar.

\section{Imagens e elementos gráficos}

Toda imagem informativa deve possuir texto alternativo ou descrição equivalente. A descrição comunica finalidade, informação essencial e relação com o argumento, sem repetir mecanicamente a legenda. Imagens decorativas devem ser marcadas como tais quando a ferramenta oferecer essa opção.

Não dependa apenas de cor. Preserve contraste, use rótulos, padrões ou formas adicionais e verifique se o conteúdo permanece compreensível em escala de cinza. Gráficos complexos podem exigir uma síntese textual dos dados mais relevantes.

\section{Tabelas}

Tabelas precisam de cabeçalhos claros e estrutura simples. Evite células mescladas quando prejudicarem a navegação, indique unidades e não use espaços ou tabulações para simular colunas. Tabelas extensas devem ser divididas de forma lógica ou acompanhadas de versão acessível complementar.

\section{Linguagem, fonte e links}

Prefira linguagem clara, frases diretas e termos técnicos explicados na primeira ocorrência. Mantenha fonte legível, tamanho suficiente e espaçamento consistente. Links devem descrever o destino; expressões como ``clique aqui'' não informam o conteúdo fora de contexto.

\section{Conferência}

Antes da entrega, tente selecionar e copiar o texto do PDF, navegue pelos títulos, examine a ordem de leitura, verifique textos alternativos e, quando possível, teste o arquivo com leitor de tela. A validação automática ajuda, mas não substitui a leitura humana.

\section{Slides e pôsteres}

Apresentações devem usar contraste adequado, pouco texto por tela, títulos informativos, ordem de leitura previsível e descrição oral das imagens relevantes. Vídeos precisam de legendas quando utilizados. Pôsteres devem preservar legibilidade à distância e não depender apenas de códigos de cor.

\fluxopendente

A UnDF deve definir requisitos mínimos, ferramentas de teste, responsabilidades e suporte para acessibilidade dos documentos, modelos, apresentações e repositório.

\chapter{Entrega e Arquivamento}

A entrega do TCC deve seguir as orientações do curso e da instituição. Além do conteúdo acadêmico, é necessário observar formato do arquivo, prazos, versão correta, identificação e eventuais documentos complementares.

\section{Versão para banca}

A versão encaminhada à banca deve estar completa, revisada e normalizada. Não é adequado enviar texto inacabado, com referências incompletas, seções sem revisão ou elementos obrigatórios ausentes.

Antes de enviar, confira:

\begin{itemize}
  \item capa e folha de rosto;
  \item resumo e palavras-chave;
  \item sumário atualizado;
  \item títulos e numeração;
  \item citações e referências;
  \item figuras, tabelas e fontes;
  \item paginação;
  \item anexos e apêndices citados;
  \item ortografia e clareza.
\end{itemize}

\section{Versão final}

A versão final é a versão corrigida após avaliação, quando houver banca ou parecer. Ela deve incorporar ajustes exigidos e passar por nova conferência formal.

O estudante deve evitar múltiplos arquivos com nomes confusos. Recomenda-se manter controle claro de versões, com data e indicação de finalidade.

\section{Formato do arquivo}

O formato de entrega deve seguir a orientação institucional. Em geral, o PDF é o formato adequado para preservar paginação, fontes, imagens e apresentação final.

Arquivos editáveis podem ser solicitados pelo curso, mas não substituem o PDF da versão destinada à entrega quando a apresentação precisa ser preservada.

\section{Arquivamento}

O estudante deve guardar cópias do trabalho final, fontes utilizadas, instrumentos, dados autorizados, pareceres, certificados e comprovantes de submissão. O arquivamento responsável facilita comprovação, revisão futura e continuidade acadêmica.

\section{Publicação e divulgação}

Quando o TCC for publicado em repositório, evento, artigo ou outro meio, devem ser observados direitos autorais, autorização de participantes, sigilo, licenças de imagem e eventuais restrições institucionais.

Resultados de pesquisa não devem ser divulgados de forma que exponha participantes ou instituições sem autorização.

\chapter{Integridade Acadêmica e Uso de Tecnologias}
\label{chap:integridade}

A integridade acadêmica sustenta a confiança na produção universitária. Ela envolve honestidade intelectual, respeito às fontes, autoria responsável, transparência metodológica e compromisso com a verdade possível dentro dos limites da pesquisa.

\section{Plágio}

Plágio ocorre quando ideias, textos, dados, imagens, estruturas ou resultados de outra pessoa são apresentados como se fossem próprios. Pode ocorrer por cópia literal sem citação, paráfrase sem referência, tradução sem indicação, uso de trabalho alheio ou reaproveitamento indevido de texto próprio sem transparência.

Mesmo quando não há intenção de enganar, descuido com citações e referências pode gerar problema grave. Por isso, o estudante deve registrar fontes desde o início da pesquisa.

\section{Tipos comuns de plágio}

\begin{longtable}{p{0.28\linewidth}p{0.60\linewidth}}
\toprule
Tipo & Caracterização \\
\midrule
\endfirsthead
\toprule
Tipo & Caracterização \\
\midrule
\endhead
Cópia literal & Reprodução de trecho sem aspas, destaque ou indicação correta da fonte. \\
Paráfrase sem fonte & Reescrita de ideia de outro autor sem citação. \\
Plágio mosaico & Combinação de trechos de várias fontes com pequenas alterações, sem autoria clara. \\
Tradução sem indicação & Tradução de texto estrangeiro apresentada como redação própria. \\
Uso indevido de imagem ou tabela & Inserção de figura, tabela, gráfico ou esquema sem fonte ou autorização quando necessária. \\
Autoplágio & Reaproveitamento de texto próprio já avaliado ou publicado sem transparência. \\
Referência falsa & Indicação de fonte inexistente, não consultada ou que não sustenta a afirmação feita. \\
Terceirização de autoria & Entrega de texto produzido por outra pessoa ou ferramenta como se fosse autoria do estudante. \\
\bottomrule
\end{longtable}

\section{Como prevenir}

Algumas práticas reduzem o risco de plágio:

\begin{itemize}
  \item registrar a fonte completa no momento da leitura;
  \item anotar páginas de trechos importantes;
  \item separar claramente citação literal, paráfrase e comentário próprio;
  \item usar aspas ou bloco próprio em transcrições literais;
  \item revisar se toda ideia externa foi citada;
  \item conferir se todas as referências correspondem a fontes realmente consultadas.
\end{itemize}

\section{Autoplágio e reaproveitamento}

O estudante pode retomar trabalhos anteriores quando isso for pertinente e autorizado, mas deve informar a origem e evitar apresentar produção já avaliada como se fosse nova. O reaproveitamento precisa ser transparente.

\section{Fabricação e manipulação de dados}

Inventar participantes, alterar respostas, omitir resultados relevantes, selecionar apenas dados convenientes ou apresentar análise que não foi realizada compromete a validade do trabalho.

Resultados inesperados ou negativos também são resultados. O TCC deve relatar o percurso com honestidade.

Fabricação é a criação de dados, participantes, fontes ou resultados inexistentes. Falsificação é a alteração, omissão ou manipulação indevida de materiais, procedimentos ou resultados. Ambas comprometem a validade acadêmica, ainda que o texto esteja formalmente bem apresentado.

\section{Uso de inteligência artificial}

Ferramentas de inteligência artificial podem apoiar revisão linguística, geração inicial de ideias, organização, tradução, síntese, programação, criação de imagens e exploração de dados, desde que seu uso respeite as regras do curso e não substitua a autoria do estudante. O uso permitido depende da finalidade, da disciplina e do regulamento aplicável.

\begin{longtable}{p{0.30\linewidth}p{0.58\linewidth}}
\toprule
Uso & Verificação indispensável \\
\midrule
Revisão e tradução & Conferir sentido, terminologia, voz autoral e alterações de conteúdo. \\
Ideias e sínteses & Retomar as fontes originais e reconstruir o argumento com compreensão própria. \\
Programação & Testar o código, licenças, segurança, resultados e adequação metodológica. \\
Imagens & Declarar a geração quando exigido e verificar direitos, vieses e integridade visual. \\
Análise de dados & Validar método, parâmetros, código, resultados e possibilidade de reprodução. \\
Busca bibliográfica & Confirmar cada obra em catálogo, base ou fonte confiável; a ferramenta não substitui a busca. \\
Referências & Nunca aceitar referência sem verificar autoria, título, veículo, data, DOI ou endereço. \\
Produção de texto & Preservar autoria intelectual, transparência e responsabilidade; texto gerado não se torna automaticamente produção do estudante. \\
\bottomrule
\end{longtable}

O estudante não deve entregar texto produzido por ferramenta como se fosse produção própria sem revisão, compreensão e responsabilidade. Também não deve inserir referências inexistentes, dados inventados ou interpretações não verificadas.

Outros riscos incluem falsa citação, reprodução de vieses, exposição de dados pessoais ou sigilosos, violação de direitos autorais, geração de imagens enganosas, perda de confidencialidade e dependência de resultados que o estudante não consegue explicar.

Quando o curso exigir declaração de uso de ferramentas, o estudante deve descrevê-la de forma transparente. O Apêndice \ref{ap:declaracao-ia} apresenta um modelo de apoio do projeto, cuja presença não cria obrigatoriedade. Ferramentas de IA não são autoras do trabalho e não assumem responsabilidade acadêmica.

\fluxopendente

A exigência da declaração, sua posição no trabalho, o fluxo de análise e eventuais consequências dependem da política institucional e das regras do curso.

\section{Similaridade textual e avaliação de integridade}

Relatórios de similaridade são instrumentos de triagem. Um percentual, isoladamente, não prova plágio e tampouco assegura integridade. A análise deve considerar a origem das coincidências, o uso de citações, a natureza do documento, as práticas da área e a presença de apropriação indevida.

\section{Direitos autorais e licenças}

Citar a fonte é indispensável, mas nem sempre basta para autorizar a reprodução integral de imagem, instrumento, escala, mapa, fotografia, software ou outro material protegido. O estudante deve observar a Lei nº 9.610/1998, licenças de uso, termos das plataformas e autorizações necessárias. Materiais em domínio público ou com licença aberta também exigem atribuição conforme as condições aplicáveis.

\section{Responsabilidade sobre a versão final}

Independentemente das ferramentas utilizadas, o estudante é responsável por todo o conteúdo entregue: afirmações, citações, referências, dados, imagens, tabelas, conclusões e normalização.

\caixaorientacao{Uma prática simples reduz riscos: mantenha uma lista de fontes consultadas, registre páginas importantes e nunca copie trechos para o rascunho sem marcar claramente que são citações.}

\chapter{Governança, Validação e Atualização}
\label{chap:governanca}

Este manual combina normas técnicas, legislação, decisões institucionais e regras de curso. Sua publicação exige governança para evitar que uma minuta técnica seja confundida com ato oficial.

\section{Instâncias a consultar}

Devem participar da validação, conforme suas competências:

\begin{itemize}
  \item Biblioteca Central da UnDF, para normalização, dados catalográficos, depósito, preservação, informações gerais do repositório e atualização normativa;
  \item Pró-Reitoria e áreas acadêmicas competentes, para regras gerais de graduação e TCC;
  \item colegiados e cursos, para compatibilidade com PPCs e regulamentos;
  \item unidade responsável pelos modelos, para publicação, atualização e suporte;
  \item instâncias de ética, para fluxos envolvendo seres humanos, animais, material biológico e dados sensíveis;
  \item área de acessibilidade, para requisitos e testes dos documentos;
  \item assessoria jurídica, quando houver dúvida sobre direitos autorais, proteção de dados, responsabilidades ou forma de aprovação.
\end{itemize}

Essas atribuições indicam participações possíveis. O ato competente deve confirmar a unidade responsável e o fluxo definitivo, sem atribuir à Biblioteca Central ou a outra unidade funções ainda não aprovadas.

\section{Etapas antes da publicação institucional}

A publicação como orientação institucional definitiva somente deve ocorrer após:

\begin{enumerate}
  \item revisão normativa;
  \item revisão da Biblioteca Central da UnDF;
  \item validação acadêmica;
  \item compatibilização com PPCs e regulamentos;
  \item sincronização com os modelos editáveis e os PDFs distribuídos;
  \item teste das validações automáticas;
  \item teste de acessibilidade;
  \item aprovação por ato formal.
\end{enumerate}

\section{Manutenção}

Cada edição deve registrar versão, data e situação. O fluxo oficial de aprovação, vigência, responsabilidades, periodicidade e arquivamento ainda precisa ser definido. Revisões extraordinárias devem ocorrer sempre que norma, legislação, regulamento ou processo institucional relevante for alterado.

\section{Conteúdo editável e elementos padronizados}

Nos modelos institucionais, o estudante deve alterar somente os campos e conteúdos destinados à autoria. Os elementos padronizados devem ser preservados conforme o modelo disponibilizado e as orientações do curso.

\chapter{Referências Normativas e Documentos de Apoio}

Este manual considera normas técnicas e documentos institucionais relacionados à apresentação de trabalhos acadêmicos. As normas devem ser consultadas diretamente quando houver dúvida específica de interpretação.

\section{Normas ABNT consideradas}

As designações a seguir foram conferidas diretamente nos exemplares normativos disponíveis ao projeto em 7 de agosto de 2026:

\begin{itemize}
  \item ABNT NBR 14724:2024, quarta edição, de 16 de dezembro de 2024, \emph{Informação e documentação --- Trabalhos acadêmicos --- Apresentação};
  \item ABNT NBR 6028:2021, segunda edição, de 18 de maio de 2021, \emph{Informação e documentação --- Resumo, resenha e recensão --- Apresentação};
  \item ABNT NBR 6027:2012, segunda edição, de 11 de dezembro de 2012, \emph{Informação e documentação --- Sumário --- Apresentação};
  \item ABNT NBR 6024:2012, \emph{Informação e documentação --- Numeração progressiva das seções de um documento --- Apresentação};
  \item ABNT NBR 10520:2023, segunda edição, de 19 de julho de 2023, \emph{Informação e documentação --- Citações em documentos --- Apresentação};
  \item ABNT NBR 6023:2025, terceira edição, de 21 de maio de 2025, \emph{Informação e documentação --- Referências --- Elaboração}.
\end{itemize}

O projeto também utiliza, conforme a modalidade e o elemento, as ABNT NBR 15287, 6022, 10719, 6034 e 15437, além das \emph{Normas de apresentação tabular} do IBGE. A edição, a designação oficial e a situação de vigência desses documentos devem ser confirmadas no catálogo autorizado antes da publicação institucional, pois os respectivos exemplares normativos não integram o acervo local auditado nesta revisão.

\section{Decisões institucionais}

Nos pontos em que a norma não define todos os detalhes de apresentação, o projeto validou decisões institucionais para preservar unidade visual e facilitar a vida acadêmica. Exemplos incluem formato de capa, apresentação da banca, modelo de data de aprovação, paginação, títulos, tratamento de listas e padrões de resumo.

Essas decisões estão registradas nos documentos de controle e incorporadas aos modelos do projeto. A publicação oficial do manual permanece sujeita ao fluxo institucional próprio, sem transformar novamente essas decisões em propostas.

\section{Documentos institucionais da UnDF}

O conteúdo institucional foi confrontado com os seguintes instrumentos:

\begin{itemize}
  \item Lei Complementar distrital nº 987, de 26 de julho de 2021;
  \item Decreto distrital nº 42.333, de 26 de julho de 2021;
  \item Estatuto da UnDF, aprovado pela Resolução nº 3/2022;
  \item Regimento Geral da UnDF, aprovado pela Portaria nº 4/2022;
  \item Resolução nº 12/2023, sobre elaboração e reestruturação de projetos pedagógicos de curso;
  \item Regimento da Graduação, aprovado pela Instrução nº 24, de 15 de abril de 2025;
  \item PPCs, regulamentos de curso e atos posteriores aplicáveis.
\end{itemize}

Os atos oficiais devem ser consultados em sua versão vigente no Sistema Integrado de Normas Jurídicas do Distrito Federal e nos canais oficiais da UnDF.

\section{Ética, animais e proteção de dados}

Também foram considerados como documentos de apoio:

\begin{itemize}
  \item Resolução CNS nº 466/2012, sobre pesquisas envolvendo seres humanos;
  \item Resolução CNS nº 510/2016, aplicável às pesquisas em Ciências Humanas e Sociais;
  \item Lei nº 11.794/2008 e normas do Concea, sobre uso científico de animais;
  \item Lei nº 13.709/2018, Lei Geral de Proteção de Dados Pessoais;
  \item Lei nº 9.610/1998, sobre direitos autorais;
  \item Guia orientativo da Autoridade Nacional de Proteção de Dados sobre tratamento de dados pessoais para fins acadêmicos e de pesquisa;
  \item orientações do Sistema CEP/Conep e das comissões institucionais competentes.
\end{itemize}

Os guias da Autoridade Nacional de Proteção de Dados possuem caráter orientativo e não substituem a legislação nem a avaliação institucional do caso concreto.

\section{Acessibilidade}

As recomendações de acessibilidade consideram o Modelo de Acessibilidade em Governo Eletrônico (eMAG), as Diretrizes de Acessibilidade para Conteúdo Web (WCAG) e boas práticas de documentos digitais. A UnDF ainda deve definir os requisitos mínimos aplicáveis aos TCCs e aos modelos institucionais.

\section{Atualização do manual}

Este manual deve ser atualizado sempre que houver:

\begin{itemize}
  \item alteração normativa relevante;
  \item nova decisão institucional;
  \item mudança em regulamento de curso;
  \item identificação de problema recorrente em bancas;
  \item necessidade de orientar novos tipos de trabalho;
  \item atualização de procedimentos de entrega ou arquivamento.
\end{itemize}

\caixaatencao{A versão preliminar deste manual deve passar por validação institucional antes de ser publicada como orientação institucional definitiva.}

\chapter{Modelos de Referências}
\label{chap:modelos-referencias}

Este capítulo reúne estruturas de conferência da ABNT NBR 6023:2025. Campos entre colchetes angulares devem ser substituídos pelos dados reais da fonte. Não mantenha marcadores no trabalho final e não complete lacunas por suposição.

\caixaatencao{Consulte a própria fonte, a página editorial, o catálogo ou a base responsável. Geradores automáticos podem ajudar, mas cada referência precisa ser conferida. Os endereços terminados em ``.exemplo'' e o DOI com ``10.xxxx'' são marcadores didáticos visíveis e não constituem hiperlinks.}

\section{Regras comuns}

As referências são alinhadas à margem esquerda, em espaço simples e separadas entre si por uma linha em branco. O recurso tipográfico usado para destacar o título deve ser uniforme. Autoria, título, edição, local, editora e data são registrados conforme o tipo documental e os dados disponíveis.

\section{Livro impresso}

\refmodelo{SOBRENOME, Prenome. \textbf{Título}: subtítulo. Número da edição, a partir da segunda. Local: Editora, ano.}

Registre responsabilidade secundária, volume, coleção ou tradução quando forem relevantes para identificar a edição consultada.

\section{Livro eletrônico}

\refmodelo{SOBRENOME, Prenome. \textbf{Título}: subtítulo. Edição. Local: Editora, ano. Formato ou suporte, quando necessário. Disponível em: \urlexemplo{https://endereco.exemplo/livro}. Acesso em: dia mês abreviado. ano.}

Não use o endereço de uma página de resultados quando houver página estável do livro. Se houver DOI, registre-o conforme aparece na fonte e mantenha os dados exigidos para o documento consultado.

\section{Capítulo de livro}

\refmodelo{SOBRENOME, Prenome do autor da parte. Título do capítulo. In: SOBRENOME, Prenome do responsável pela obra (org.). \textbf{Título do livro}: subtítulo. Edição. Local: Editora, ano. páginas inicial--final.}

A autoria do capítulo e a responsabilidade pelo livro são informações diferentes. Não atribua o capítulo ao organizador quando houver autoria própria.

\section{Artigo de periódico sem DOI}

\refmodelo{SOBRENOME, Prenome. Título do artigo. \textbf{Título do periódico}, local, volume, número, páginas ou identificador eletrônico, período e ano. Disponível em: \urlexemplo{https://endereco.exemplo/artigo}. Acesso em: dia mês abreviado. ano.}

Em artigo impresso, omitem-se endereço e acesso. Use o identificador eletrônico no lugar das páginas quando essa for a forma adotada pelo periódico.

\section{Artigo de periódico com DOI}

\refmodelo{SOBRENOME, Prenome. Título do artigo. \textbf{Título do periódico}, local, volume, número, páginas ou identificador, ano. DOI: \urlexemplo{https://doi.org/10.xxxx/xxxxx}. Disponível em: \urlexemplo{https://endereco-consultado.exemplo}. Acesso em: dia mês abreviado. ano.}

O DOI é um identificador persistente; a URL informa onde o conteúdo foi consultado. Não invente DOI e não use a expressão ``DOI ou Disponível em'' como se os elementos fossem sempre intercambiáveis. Quando a consulta for eletrônica, registre os elementos exigidos para acesso conforme a norma e o tipo documental.

\section{TCC, dissertação e tese}

\refmodelo{SOBRENOME, Prenome. \textbf{Título}: subtítulo. Ano de depósito. Tipo do trabalho (Grau e área) -- Unidade acadêmica, Instituição, Local, ano de defesa.}

Em versão eletrônica, acrescente o endereço permanente e a data de acesso. Diferencie o ano de depósito do ano de defesa quando a fonte apresentar ambos.

\section{Trabalho em evento}

\refmodelo{SOBRENOME, Prenome. Título do trabalho. In: NOME DO EVENTO, número., ano, local de realização. \textbf{Título dos anais [...]} Local de publicação: Responsável pela publicação, ano. páginas ou identificador. Disponível em: \urlexemplo{https://endereco.exemplo}. Acesso em: dia mês abreviado. ano.}

Não confunda local do evento com local de publicação. Para evento on-line, registre os dados efetivamente apresentados pela fonte.

\section{Legislação}

\refmodelo{BRASIL. Lei nº 13.709, de 14 de agosto de 2018. Lei Geral de Proteção de Dados Pessoais. \textbf{Diário Oficial da União}: Brasília, DF, 15 ago. 2018. Disponível em: \url{https://www.planalto.gov.br/ccivil_03/_ato2015-2018/2018/lei/l13709.htm}. Acesso em: 6 ago. 2026.}

Identifique a jurisdição, o tipo e número do ato, a data, a ementa ou título, a publicação oficial e os dados de acesso. Consulte o texto consolidado quando houver alterações.

\section{Jurisprudência}

\refmodelo{JURISDIÇÃO. Órgão julgador. Tipo e número do processo. Relator: [nome]. Data do julgamento. \textbf{Fonte de publicação}, local, data. Disponível em: \urlexemplo{https://endereco-oficial.exemplo}. Acesso em: dia mês abreviado. ano.}

Registre os elementos capazes de identificar a decisão específica. Não referencie apenas a página inicial do tribunal.

\section{Ato administrativo}

\refmodelo{DISTRITO FEDERAL. Universidade do Distrito Federal Professor Jorge Amaury Maia Nunes. Instrução nº 24, de 15 de abril de 2025. Aprova o Regimento da Graduação da UnDF. \textbf{Diário Oficial do Distrito Federal}: Brasília, DF, nº 74, p. 10, 22 abr. 2025.}

Portarias, instruções, resoluções e editais devem identificar a entidade, o ato, a data, a ementa e a fonte oficial. Acrescente endereço e acesso quando consultados on-line.

\section{Documento institucional}

\refmodelo{ENTIDADE RESPONSÁVEL. \textbf{Título do documento}: subtítulo. Versão ou edição. Local: Entidade, ano. Disponível em: \urlexemplo{https://endereco-institucional.exemplo}. Acesso em: dia mês abreviado. ano.}

Use a instituição como autora somente quando ela assumir responsabilidade intelectual pelo documento.

\section{Relatório técnico ou científico}

\refmodelo{AUTOR OU ENTIDADE. \textbf{Título do relatório}: subtítulo. Local: Instituição responsável, ano. Número ou série do relatório. Disponível em: \urlexemplo{https://endereco.exemplo}. Acesso em: dia mês abreviado. ano.}

Inclua número, contrato, projeto ou série quando forem necessários para distinguir o relatório.

\section{Norma técnica}

\refmodelo{ASSOCIAÇÃO BRASILEIRA DE NORMAS TÉCNICAS. \textbf{ABNT NBR 10520: informação e documentação -- citações em documentos -- apresentação}. Rio de Janeiro: ABNT, 2023.}

Use o código, o título e o ano da edição efetivamente consultada. Não cite como vigente uma edição substituída.

\section{Matéria jornalística}

\refmodelo{SOBRENOME, Prenome. Título da matéria. \textbf{Nome do jornal}, local, dia mês abreviado. ano. Seção, página.}

Na versão eletrônica, acrescente endereço e acesso. Quando não houver autoria pessoal, a entrada pode começar pelo título, conforme a regra aplicável.

\section{Vídeo}

\refmodelo{TÍTULO do vídeo. Responsável principal: [nome]. [Local]: Produtora ou canal, ano. 1 vídeo (duração). Disponível em: \urlexemplo{https://endereco.exemplo}. Acesso em: dia mês abreviado. ano.}

Registre a responsabilidade relevante, a plataforma ou produtora, a data e a duração. Use timestamp na citação direta de trecho audiovisual.

\section{Podcast}

\refmodelo{TÍTULO do episódio. In: \textbf{Título do podcast}. Apresentação: [nome]. [Local]: Produtora, dia mês abreviado. ano. Podcast (duração). Disponível em: \urlexemplo{https://endereco.exemplo}. Acesso em: dia mês abreviado. ano.}

Distinga o título do episódio do título da série e identifique apresentador, entrevistado ou produtor conforme a responsabilidade efetiva.

\section{Entrevista}

\refmodelo{SOBRENOME, Prenome do entrevistado. \textbf{Título ou assunto}. Entrevista concedida a [nome do entrevistador]. Local, data. Suporte, duração ou localização, quando necessário.}

Entrevistas produzidas pela pesquisa também precisam ser descritas metodologicamente. Use código do participante quando a identificação for vedada e preserve os registros segundo o protocolo ético.

\section{Imagem, fotografia ou ilustração}

\refmodelo{AUTOR OU ENTIDADE. \textbf{Título}. Ano. 1 imagem, técnica ou tipo, dimensões quando relevantes. Acervo ou plataforma. Disponível em: \urlexemplo{https://endereco.exemplo}. Acesso em: dia mês abreviado. ano.}

Além da referência, verifique autoria, licença e autorização de uso. ``Fonte: Google'' não identifica a obra nem concede licença.

\section{Mapa ou documento cartográfico}

\refmodelo{AUTOR OU ENTIDADE. \textbf{Título}: subtítulo. Local: Responsável, ano. Escala. Tipo e suporte. Disponível em: \urlexemplo{https://endereco.exemplo}. Acesso em: dia mês abreviado. ano.}

Registre escala, projeção, folha, série e coordenadas quando forem relevantes para identificar ou reproduzir a análise.

\section{Página da internet}

\refmodelo{AUTOR OU ENTIDADE. \textbf{Título da página}. Local ou portal, data de publicação ou atualização. Disponível em: \urlexemplo{https://endereco-especifico.exemplo}. Acesso em: dia mês abreviado. ano.}

Prefira a página específica. Não use apenas o nome do buscador, a página inicial do portal ou um endereço encurtado se houver URL estável.

\section{Base de dados}

\refmodelo{ENTIDADE RESPONSÁVEL. \textbf{Nome da base de dados}. Versão ou período. Local: Entidade, ano. Disponível em: \urlexemplo{https://endereco.exemplo}. Acesso em: dia mês abreviado. ano.}

Informe versão, recorte, tabela, painel, filtros e data de extração no texto ou na metodologia quando os resultados puderem mudar.

\section{Conjunto de dados}

\refmodelo{AUTOR OU ENTIDADE. \textbf{Título do conjunto de dados}. Versão. Repositório, ano. Conjunto de dados. DOI: \urlexemplo{https://doi.org/10.xxxx/xxxxx}. Disponível em: \urlexemplo{https://endereco.exemplo}. Acesso em: dia mês abreviado. ano.}

Indique versão, responsável, repositório, identificador persistente e data de acesso. Não confunda a referência do conjunto com a do artigo que o analisou.

\section{Publicação em rede social}

\refmodelo{AUTOR OU PERFIL RESPONSÁVEL. Início ou título da publicação. Local: Plataforma, dia mês abreviado. ano. Tipo de publicação. Disponível em: \urlexemplo{https://endereco-direto.exemplo}. Acesso em: dia mês abreviado. ano.}

Use redes sociais somente quando forem fonte pertinente ao objeto. Preserve evidência da publicação quando a mutabilidade for relevante e observe ética, privacidade e termos de uso.

\section{Documento sem autoria pessoal}

Procure primeiro autoria institucional. Se nenhuma responsabilidade puder ser identificada, a entrada começa pelo título conforme a NBR 6023:2025. A ausência de autoria exige avaliação crítica reforçada e não autoriza atribuir o conteúdo a uma instituição apenas porque ela hospeda o arquivo.

\section{Local, editora ou data não identificados}

Não omita silenciosamente elementos nem invente dados. Use informação inferida entre colchetes quando houver base segura. Se o local não puder ser identificado, empregue [S. l.]; se a editora não puder ser identificada, [s. n.]; se ambos faltarem, [S. l.: s. n.]. Para data ausente, indique uma data provável ou intervalo entre colchetes, como [2024?] ou [entre 2020 e 2024], somente quando a inferência puder ser justificada.

\section{Conferência final}

Antes da entrega, confira:

\begin{itemize}
  \item autoria e responsabilidade conforme a fonte;
  \item título, edição, local, editora e data;
  \item DOI, URL estável e data de acesso quando aplicáveis;
  \item funcionamento dos links e correspondência com o documento citado;
  \item uniformidade de destaque, pontuação e abreviaturas;
  \item presença de toda fonte citada na lista;
  \item uso efetivo de toda fonte incluída na lista.
\end{itemize}

\chapter{Problemas Recorrentes e Como Evitá-los}

Este capítulo reúne problemas frequentes em TCCs e orientações para preveni-los. A intenção é ajudar o estudante a revisar o trabalho antes que essas falhas apareçam na banca.

\section{Tema amplo demais}

Um tema amplo torna a pesquisa genérica. O estudante tenta falar de tudo, mas não aprofunda nada. Para corrigir, delimite espaço, tempo, população, objeto, material ou enfoque.

Pergunte: qual parte desse tema eu consigo investigar com rigor no tempo disponível?

\section{Problema que não é pergunta de pesquisa}

Muitos trabalhos apresentam problema como assunto, opinião ou justificativa. O problema precisa orientar investigação. Se ele não pode ser respondido pelo método, ainda não está pronto.

\section{Objetivos desconectados}

Os objetivos específicos devem conduzir ao objetivo geral. Quando cada objetivo aponta para direção diferente, o trabalho perde unidade.

Revise se cada objetivo específico aparece depois na metodologia, nos resultados e na conclusão.

\section{Revisão de literatura descritiva}

Uma revisão frágil resume autor por autor, sem articulação. A revisão deve construir debate. Agrupe autores por ideias, conceitos, convergências e divergências.

\section{Metodologia genérica}

Frases como ``foi realizada pesquisa qualitativa'' ou ``foi feita revisão bibliográfica'' são insuficientes. A metodologia precisa explicar procedimentos concretos.

\section{Resultados sem discussão}

Apresentar dados não é discutir dados. A discussão interpreta resultados, relaciona com objetivos e dialoga com literatura.

\section{Conclusão sem resposta}

A conclusão deve responder ao problema. Se ela apenas repete a introdução, falta fechamento analítico.

\section{Citações em excesso}

Excesso de citações pode apagar a voz do estudante. Use fontes para sustentar o argumento, mas conduza o texto com análise própria.

\section{Referências inconsistentes}

Referências incompletas passam impressão de descuido e dificultam verificação. Faça conferência cruzada: texto para lista final e lista final para texto.

\section{Elementos visuais sem função}

Figuras e tabelas precisam ser necessárias. Se o texto não comenta o elemento, talvez ele não deva estar ali.

\section{Normalização deixada para o fim}

Normalizar apenas na última etapa pode revelar problemas difíceis de corrigir. Ainda assim, a revisão final deve ser feita no PDF, porque quebras de página e listas só aparecem plenamente na versão final.

\section{Falta de revisão após correções}

Alterações feitas depois da banca podem criar erros novos. Sempre revise sumário, paginação, referências e coerência depois de qualquer mudança significativa.

\caixaorientacao{A melhor prevenção é trabalhar por versões. Cada versão deve resolver um tipo de problema: estrutura, conteúdo, método, citações, linguagem e apresentação.}



\setcounter{capitulosparteum}{\value{chapter}}
\appendix
\chapter{Roteiro para Projeto de TCC}

Este roteiro pode ser utilizado como apoio para elaboração do projeto de TCC. O curso pode adaptar, reduzir ou ampliar os itens conforme sua área.

O roteiro segue a lógica da ABNT NBR 15287:2025, mas não substitui o formulário ou as exigências aprovadas pelo curso.

\section{Elementos pré-textuais}

Utilize folha de rosto com autoria, título, natureza do projeto, vínculo acadêmico, orientação, local e ano. Inclua capa, listas e sumário quando forem exigidos ou necessários. Não antecipe no projeto elementos próprios da versão final, como folha de aprovação.

\section{Identificação}

\begin{itemize}
  \item Nome do estudante.
  \item Curso.
  \item Orientador.
  \item Título provisório.
  \item Linha de pesquisa ou área, quando houver.
\end{itemize}

\section{Tema e delimitação}

Descreva o tema geral e o recorte adotado. Indique espaço, tempo, população, material, objeto ou enfoque. A delimitação deve mostrar o que será estudado e o que ficará fora do trabalho.

\section{Problema de pesquisa}

Apresente a pergunta central do trabalho. Ela deve ser clara, investigável e compatível com o método.

\section{Objetivos}

Apresente um objetivo geral e objetivos específicos. Os objetivos específicos devem indicar etapas intelectuais necessárias para alcançar o objetivo geral.

\section{Justificativa}

Explique a relevância do trabalho. Considere importância acadêmica, social, profissional, institucional ou tecnológica.

\section{Referencial teórico inicial}

Apresente autores, conceitos e debates que sustentam a pesquisa. Mostre como essas fontes ajudam a compreender o problema.

\section{Metodologia}

Indique tipo de pesquisa, abordagem, materiais, participantes, instrumentos, procedimentos, critérios de análise e cuidados éticos.

\section{Recursos e riscos}

Registre recursos necessários, formas de acesso e riscos capazes de afetar a execução. Indique medidas alternativas para autorizações, recrutamento, coleta, equipamentos ou bases essenciais.

\section{Cronograma}

\begin{longtable}{p{0.45\linewidth}p{0.37\linewidth}}
\toprule
Atividade & Período previsto \\
\midrule
\endfirsthead
\toprule
Atividade & Período previsto \\
\midrule
\endhead
Delimitação do tema e problema & \\
Levantamento bibliográfico & \\
Elaboração do projeto & \\
Submissão ética, quando aplicável & \\
Coleta ou seleção de dados & \\
Análise & \\
Redação das seções textuais & \\
Revisão com orientador & \\
Normalização & \\
Entrega para banca & \\
Correções finais & \\
\bottomrule
\end{longtable}

\section{Referências iniciais}

Liste as fontes que sustentam a proposta. O projeto não precisa conter bibliografia definitiva, mas deve demonstrar base mínima para início da pesquisa.

\section{Apêndices, anexos e índice}

Inclua instrumentos elaborados pela autoria como apêndices e documentos de terceiros como anexos, somente quando necessários. O índice é opcional e deve seguir a ABNT NBR 6034:2004 quando utilizado.

\chapter{Matriz de Coerência do TCC}

A matriz de coerência ajuda a verificar se as partes do trabalho estão alinhadas. Ela pode ser usada no projeto, durante a escrita e antes da entrega.

\begin{longtable}{p{0.26\linewidth}p{0.62\linewidth}}
\toprule
Item & Pergunta de verificação \\
\midrule
\endfirsthead
\toprule
Item & Pergunta de verificação \\
\midrule
\endhead
Tema & O tema está delimitado e é compatível com o curso? \\
Problema & A pergunta de pesquisa é clara e investigável? \\
Objetivo geral & O objetivo geral responde diretamente ao problema? \\
Objetivos específicos & Os objetivos específicos conduzem ao objetivo geral? \\
Justificativa & A relevância do trabalho está demonstrada com argumentos e fontes? \\
Referencial teórico & Os conceitos utilizados ajudam a analisar o problema? \\
Metodologia & O método permite responder ao problema? \\
Dados ou corpus & Os materiais analisados são adequados e suficientes? \\
Resultados & Os resultados se relacionam aos objetivos? \\
Discussão & A análise dialoga com a literatura? \\
Conclusão & A conclusão responde ao problema e reconhece limites? \\
Referências & Todas as fontes citadas estão corretamente registradas? \\
\bottomrule
\end{longtable}

\section{Como usar}

Preencha a matriz em linguagem simples. Se algum item não puder ser respondido com segurança, o trabalho provavelmente precisa de ajuste. A matriz também pode ser apresentada ao orientador para facilitar a discussão sobre coerência do projeto.

\caixaorientacao{A matriz não é parte obrigatória do TCC final. Ela é um instrumento de planejamento e revisão.}


\chapter{Checklists Acadêmicos}

Os checklists abaixo apoiam a revisão por etapas. Eles não substituem a leitura crítica, a orientação nem os critérios específicos do curso.

\section{Checklist do projeto}

\begin{itemize}
  \item O tema está delimitado e o problema é investigável.
  \item O objetivo geral responde ao problema e os objetivos específicos descrevem etapas necessárias.
  \item A justificativa demonstra relevância acadêmica, social, profissional ou institucional.
  \item O referencial inicial reúne fontes pertinentes e atuais para o objeto.
  \item A metodologia é compatível com os objetivos e descreve análise, limitações e gestão dos dados.
  \item O cronograma é realista e reserva tempo para revisão.
  \item Riscos, autorizações e necessidade de avaliação ética foram verificados antes da coleta.
  \item As referências utilizadas no projeto estão registradas.
\end{itemize}

\section{Checklist para orientação}

\begin{itemize}
  \item A versão enviada está identificada e permite localizar as alterações.
  \item As solicitações anteriores foram atendidas ou justificadas.
  \item As dúvidas para a reunião são objetivas.
  \item Problema, objetivos, método, análise e parte conclusiva permanecem coerentes.
  \item Citações e referências novas estão registradas.
  \item Lacunas, limitações e seções ainda incompletas estão sinalizadas.
  \item O estudante registrou as decisões tomadas durante a orientação.
\end{itemize}

\section{Checklist para banca}

\begin{itemize}
  \item O trabalho possui os elementos obrigatórios aplicáveis.
  \item Resumos, palavras-chave, sumário e listas foram conferidos.
  \item Objetivos, resultados e parte conclusiva são coerentes.
  \item A metodologia permite compreender o percurso realizado.
  \item Figuras, tabelas, quadros e demais objetos são citados e possuem título e fonte.
  \item Citações e referências foram conferidas em ambos os sentidos.
  \item Apêndices e anexos são pertinentes e foram mencionados no texto.
  \item A paginação e a legibilidade foram verificadas no PDF enviado.
  \item O envio foi autorizado conforme a regra do curso.
\end{itemize}

\section{Checklist da versão corrigida}

\begin{itemize}
  \item As correções da banca foram incorporadas ou justificadas à instância competente.
  \item Alterações de conteúdo não criaram contradições em outras seções.
  \item Sumário, listas, paginação e referências foram atualizados.
  \item O PDF abre corretamente e corresponde ao arquivo autorizado para entrega.
  \item Nome, formato, prazo e canal seguem a orientação institucional vigente.
  \item Restrições, embargo, termos ou autorizações foram anexados quando aplicáveis.
\end{itemize}

\section{Checklist de normalização}

\begin{itemize}
  \item Papel, margens, fonte, corpo, espaçamento e parágrafos estão uniformes.
  \item A contagem e a exibição da paginação estão corretas.
  \item A hierarquia de seções é progressiva e não ultrapassa o nível quinário.
  \item Títulos do texto e do sumário coincidem.
  \item Siglas, abreviaturas, alíneas e subalíneas seguem um padrão estável.
  \item Citações diretas possuem localização quando aplicável.
  \item Referências apresentam dados suficientes e não contêm marcadores não preenchidos.
  \item Elementos pré-textuais não aparecem no sumário.
\end{itemize}

\section{Checklist de acessibilidade}

\begin{itemize}
  \item O PDF é pesquisável e possui título e idioma definidos nos metadados.
  \item A ordem de leitura e a hierarquia de títulos são compreensíveis.
  \item Imagens informativas possuem texto alternativo ou descrição equivalente.
  \item Contraste, tamanho do texto e espaçamento favorecem a leitura.
  \item Informações não dependem exclusivamente de cor.
  \item Tabelas possuem cabeçalhos claros e estrutura simples.
  \item Links têm texto descritivo e funcionam no arquivo entregue.
  \item Slides e pôsteres, quando utilizados, seguem os mesmos princípios.
\end{itemize}

\section{Checklist de integridade}

\begin{itemize}
  \item Ideias, textos, dados, imagens e códigos de terceiros estão atribuídos.
  \item Paráfrases foram realmente reescritas e continuam acompanhadas de fonte.
  \item Não há referências inventadas, dados fabricados ou resultados manipulados.
  \item Reutilização de produção anterior foi identificada e autorizada quando necessário.
  \item Contribuições de outras pessoas não foram apresentadas como autoria própria.
  \item O uso de inteligência artificial, quando houver, respeita as regras aplicáveis e foi declarado quando exigido.
  \item Informações pessoais, sigilosas ou protegidas não foram expostas indevidamente.
  \item O estudante revisou o conteúdo e assume responsabilidade pela versão entregue.
\end{itemize}

\chapter{Rubrica Sugerida de Avaliação}

Esta rubrica é uma sugestão de apoio. Cada curso pode definir critérios, pesos e instrumentos próprios. Use-a também de forma formativa para identificar fragilidades e discutir prioridades, sem transformar o julgamento acadêmico em checklist mecânico.

\begin{longtable}{p{0.22\linewidth}p{0.30\linewidth}p{0.30\linewidth}}
\toprule
Critério & Desempenho satisfatório & Pontos de atenção \\
\midrule
\endfirsthead
\toprule
Critério & Desempenho satisfatório & Pontos de atenção \\
\midrule
\endhead
Tema e problema & Tema delimitado, problema claro e relevante. & Tema amplo, pergunta vaga ou recorte inviável. \\
Objetivos & Objetivos coerentes e verificáveis. & Objetivos genéricos ou desconectados do método. \\
Justificativa & Relevância demonstrada com argumentos e fontes. & Justificativa opinativa ou superficial. \\
Referencial teórico & Conceitos articulados ao problema. & Revisão descritiva, desatualizada ou desconexa. \\
Metodologia & Procedimentos claros e adequados. & Método genérico, incompleto ou incompatível. \\
Análise & Resultados interpretados com rigor. & Dados apenas descritos ou sem diálogo teórico. \\
Conclusão & Responde ao problema e reconhece limites. & Repete introdução ou traz informações novas. \\
Escrita & Texto claro, coeso e acadêmico. & Ambiguidade, informalidade ou problemas recorrentes de redação. \\
Citações e referências & Fontes adequadas e padronizadas. & Ausências, inconsistências ou uso inadequado de citações. \\
Apresentação formal & Trabalho conforme padrão institucional. & Paginação, sumário, listas ou elementos obrigatórios com falhas. \\
Defesa oral & Domínio do percurso e respostas consistentes. & Desconhecimento do próprio texto ou respostas improvisadas. \\
\bottomrule
\end{longtable}

\chapter{Modelo de Cronograma de TCC}

O cronograma abaixo é apenas um modelo. Cada curso e cada estudante devem adaptá-lo ao calendário acadêmico, à modalidade de TCC e às exigências de orientação.

\begin{longtable}{p{0.36\linewidth}p{0.12\linewidth}p{0.12\linewidth}p{0.12\linewidth}p{0.12\linewidth}}
\toprule
Atividade & Mês 1 & Mês 2 & Mês 3 & Mês 4 \\
\midrule
\endfirsthead
\toprule
Atividade & Mês 1 & Mês 2 & Mês 3 & Mês 4 \\
\midrule
\endhead
Definição do tema & & & & \\
Delimitação e problema & & & & \\
Objetivos e justificativa & & & & \\
Levantamento bibliográfico & & & & \\
Projeto ou plano de trabalho & & & & \\
Consulta ética e autorizações & & & & \\
Submissão ao sistema competente, quando aplicável & & & & \\
Coleta ou seleção de dados & & & & \\
Análise dos dados & & & & \\
Redação das seções textuais & & & & \\
Revisão com orientador & & & & \\
Normalização & & & & \\
Entrega para banca & & & & \\
Correções finais & & & & \\
\bottomrule
\end{longtable}

\section{Orientações de uso}

O cronograma deve ser construído de trás para frente a partir da data final de entrega. Reserve tempo para atrasos, revisão e ajustes. Se houver comitê de ética, autorização institucional ou coleta externa, amplie o prazo.

Não inicie recrutamento ou coleta antes das autorizações e aprovações aplicáveis. Como esses prazos podem variar, registre dependências e alternativas em vez de concentrar todas as etapas em um único mês.

\chapter{Modelo de Declaração de Uso de Inteligência Artificial}
\label{ap:declaracao-ia}

Este é um modelo de apoio do projeto. Sua inclusão no manual não torna a declaração obrigatória: o curso pode exigir informações adicionais, adotar outro formulário ou proibir determinados usos.

\modeloinstitucional{DECLARAÇÃO DE USO DE INTELIGÊNCIA ARTIFICIAL}{%
Declaro que utilizei a ferramenta [nome e versão, quando disponível] na etapa [etapa do trabalho], com a finalidade de [finalidade]. O conteúdo processado consistiu em [descrição sem expor dados pessoais, sigilosos ou protegidos].\\[0.8em]

O resultado foi utilizado para [forma de uso] e submetido à revisão humana por meio de [procedimentos de verificação]. As fontes, citações, referências, dados, cálculos, código, imagens e conclusões foram conferidos sob minha responsabilidade.\\[0.8em]

Declaro que a ferramenta não foi indicada como autora e que assumo integral responsabilidade pelo conteúdo final do trabalho.\\[1.5em]

Brasília, DF, [dia] de [mês] de [ano].\\[1.5em]

Nome e assinatura do estudante.}

\section{Quando declarar}

A exigência, a posição no trabalho e o nível de detalhamento dependem de decisão institucional e das regras do curso. Mesmo quando não houver formulário obrigatório, o uso metodologicamente relevante deve ser descrito na seção adequada do trabalho.

\chapter{Exemplo de Contagem e Exibição da Paginação}
\label{ap:paginacao}

O quadro mostra um trabalho hipotético. A quantidade real de páginas pré-textuais varia conforme os elementos utilizados.

\begin{longtable}{p{0.34\linewidth}p{0.20\linewidth}p{0.18\linewidth}p{0.16\linewidth}}
\toprule
Elemento & Participa da contagem? & Número lógico & Número visível? \\
\midrule
Capa & Não & -- & Não \\
Folha de rosto & Sim & 1 & Não \\
Dados catalográficos & Não & -- & Não \\
Folha de aprovação & Sim & 2 & Não \\
Dedicatória & Sim & 3 & Não \\
Agradecimentos & Sim & 4 & Não \\
Resumo & Sim & 5 & Não \\
Resumo estrangeiro & Sim & 6 & Não \\
Lista de ilustrações & Sim & 7 & Não \\
Sumário & Sim & 8 & Não \\
Primeira página da introdução & Sim & 9 & Sim: 9 \\
Demais textuais e pós-textuais & Sim & Sequencial & Sim \\
\bottomrule
\end{longtable}

\caixaatencao{Não se imprime numeração romana nos elementos pré-textuais do TCC. A tabela usa “número lógico” apenas para explicar a contagem que será exibida a partir da primeira página textual.}


\renewcommand{\thechapter}{\arabic{chapter}}
\renewcommand{\theHchapter}{parteii.\arabic{chapter}}
\makeatletter
\gdef\@chapapp{Capítulo}
\makeatother
\setcounter{chapter}{\value{capitulosparteum}}
\setcounter{section}{0}

\part{Guia de Uso do \marcaUnDFTeX}

\chapter*{Apresentação da Parte II}
\addcontentsline{toc}{chapter}{Apresentação da Parte II}

Esta parte documenta o uso do UnDF-TeX, desde a preparação do ambiente até a geração e a conferência do PDF. Seu público inclui estudantes que adotam o template, pessoas que prestam apoio técnico e responsáveis por manutenção.

O guia explica implementação e operação. As regras acadêmicas, metodológicas e de normalização pertencem à Parte I. Quando um comando produz capa, citação, tabela ou outro elemento, esta parte descreve como usar o recurso; a Parte I explica o que o elemento significa e quando deve ser utilizado.

\manualwarning{O UnDF-TeX 1.0.0 é um projeto técnico em validação. Sua existência não torna o uso obrigatório e não substitui PPC, regulamento de curso, orientação, decisão de colegiado ou ato institucional.}

\chapter{Apresentação da Ferramenta}

\section{O que é o UnDF-TeX}

O UnDF-TeX é uma classe e um conjunto de arquivos LaTeX destinados à preparação de Trabalhos de Conclusão de Curso. O projeto separa conteúdo, metadados, elementos opcionais, bibliografia e implementação interna para reduzir formatação manual e tornar a saída reproduzível.

O template gera capa, folha de rosto, folha de aprovação, resumos, sumário, capítulos, referências e elementos opcionais a partir dos dados e arquivos preenchidos pelo usuário. Também oferece comandos para citações, figuras, tabelas, fluxogramas, equações, pseudocódigos e referências cruzadas.

\section{Finalidade e limites}

O objetivo é automatizar tarefas repetitivas e detectar problemas técnicos comuns. O template não avalia mérito acadêmico, não verifica a veracidade das fontes, não decide se uma pesquisa precisa de autorização ética e não substitui revisão humana.

\manualnote{Automação de aparência não transforma uma escolha do projeto em regra da ABNT nem em decisão institucional. O manual acadêmico e as orientações do curso continuam sendo as referências para conteúdo e aplicabilidade.}

\section{Versões e situação}

Este guia corresponde à revisão documental \texttt{1.0.0} e descreve o UnDF-TeX \texttt{1.0.0}. Antes de iniciar ou atualizar um trabalho, confira o arquivo \texttt{VERSION.txt}, o \href{https://github.com/igormrqs/undftex-publico/blob/main/CHANGELOG.md}{histórico de alterações} e a origem do pacote recebido. A versão do manual e a versão técnica do template têm controles próprios.

Versões do projeto podem mudar. Preserve uma cópia funcional do trabalho antes de substituir a classe ou os módulos internos e faça nova regressão do PDF depois da atualização.

\section{Relação com o manual acadêmico}

Este guia remete ao Manual de TCC sempre que o usuário precisa decidir se um elemento é obrigatório, opcional ou adequado. Aqui são documentados apenas localização, sintaxe, processamento, mensagens e comportamento técnico.

\chapter{Preparação do Ambiente}

\section{Requisitos}

O projeto exige uma distribuição LaTeX atualizada com LuaLaTeX e Biber. Também é necessário um editor de texto com suporte a arquivos UTF-8 e, para conferência, um leitor de PDF.

Ambientes locais podem usar MiKTeX ou TeX Live. Em serviços online, selecione LuaLaTeX como compilador e envie a estrutura completa do projeto, preservando nomes e caminhos relativos.

\section{Obtenção da versão correta}

Obtenha o pacote pelo repositório controlado do projeto ou por canal institucional indicado. Não copie apenas \texttt{tcc.tex}: a classe, os módulos de \texttt{src/}, a identidade visual e os arquivos bibliográficos são necessários.

Depois de extrair ou clonar o projeto, confirme a presença de:

\begin{lstlisting}
undftex.cls
tcc.tex
src/
identidade/
obrigatorio/
opcional/
capitulos/
\end{lstlisting}

\section{Teste inicial}

Compile o modelo em branco antes de inserir conteúdo. O teste separa problemas de instalação de erros introduzidos durante a edição.

\begin{lstlisting}
lualatex -interaction=nonstopmode -halt-on-error tcc.tex
biber tcc
lualatex -interaction=nonstopmode -halt-on-error tcc.tex
lualatex -interaction=nonstopmode -halt-on-error tcc.tex
\end{lstlisting}

O resultado esperado é \texttt{tcc.pdf}. Mensagens sobre campos ainda não preenchidos são normais no primeiro teste; interrupção do processo, arquivo ausente ou comando desconhecido exigem correção.

\section{Editor e codificação}

Mantenha os arquivos em UTF-8. Não substitua caracteres acentuados por sequências incompatíveis nem altere a codificação de apenas um capítulo. Configure o editor para mostrar final de linha, pares de chaves e erros de sintaxe.

\chapter{Organização dos Arquivos}

\section{Entradas públicas}

A entrada da monografia é \texttt{tcc.tex}. Para consultar recursos demonstrativos,
use \texttt{tcc-exemplo-undf.tex} e a pasta \texttt{exemplo/}. Não desenvolva o
trabalho diretamente sobre o exemplo completo, pois ele contém elementos
opcionais e textos destinados a teste.

Na branch de desenvolvimento, as demais entradas são \texttt{tcc-artigo.tex}
(TCC-artigo), \texttt{artigo.tex} (artigo autônomo),
\texttt{relatorio-projeto-aplicado.tex} (relatório parcial de cada projeto),
\texttt{portfolio.tex} (relatório final do portfólio, modalidade 1.1),
\texttt{relatorio-projeto.tex}
(relatório final de projeto técnico/profissional, modalidade 1.3) e
\texttt{negocio-tecnologico.tex} (relatório final de negócio tecnológico,
modalidade 1.4). O índice \texttt{modelos/README.md} relaciona cada modelo ao
seu exemplo e aos arquivos editáveis. A base pública possui desenvolvimento
independente e não inclui automaticamente essas modalidades experimentais.

O relatório parcial preserva memória, conteúdo e evidências de cada projeto.
Sua abertura é enxuta, sem capa de TCC, aprovação, resumos ou sumário; não
exige periodicidade, parecer ou cronograma. Edite a identificação, o relato
e as fontes consultadas em \texttt{modelos/projeto-aplicado/}.

O portfólio final reúne todos os projetos do percurso declarado: introdução,
organização, inventário, síntese de cada projeto, análise integradora e
conclusão. Não basta concatenar os registros. Edite os metadados, resumos,
conteúdo e bibliografia em \texttt{modelos/portfolio/}. O exemplo preenchido
está em \texttt{portfolio-exemplo-undf.tex}, com arquivos em
\texttt{exemplo/portfolio/}. O número de projetos do exemplo é demonstrativo.

Os elementos obrigatórios e opcionais e as referências do portfólio final
preservam o padrão do TCC convencional. Artigo científico não é exigido
nesses modelos nem no relatório profissional. O documento de agosto de 2026
foi referência de organização do relato conclusivo. A admissibilidade e o
recorte do percurso devem ser confirmados com o curso.

\section{Mapa do projeto}

\begin{longtable}{p{0.31\linewidth}p{0.57\linewidth}}
\toprule
Caminho & Finalidade \\
\midrule
\endfirsthead
\toprule
Caminho & Finalidade \\
\midrule
\endhead
\texttt{tcc.tex} & Entrada do modelo em branco e ordem de montagem do documento. \\
\texttt{capitulos/} & Conteúdo textual da monografia. \\
\texttt{obrigatorio/} & Resumos e base bibliográfica da monografia. \\
\texttt{artigo/} & Metadados, resumos e miolo dos modelos de artigo. \\
\texttt{modelos/} & Índice e arquivos preenchíveis de relatório parcial, portfólio, projeto profissional e negócio tecnológico; apresentação comum dos TCCs finais. \\
\texttt{exemplo/} & Conteúdo didático independente dos modelos para preenchimento. \\
\texttt{opcional/} & Errata, dedicatória, agradecimentos, epígrafe, glossário, apêndices e anexos. \\
\texttt{figuras/} & Imagens usadas pelo trabalho. \\
\texttt{identidade/} & Ativos institucionais carregados pelo template. \\
\texttt{src/core/} & Metadados, estrutura, títulos, layout, arquivos e referências cruzadas. \\
\texttt{src/abnt/} & Camadas técnicas organizadas por norma. \\
\texttt{src/identity/} & Capa e páginas de identificação. \\
\texttt{src/biblatex/} & Estilo e comandos bibliográficos. \\
\bottomrule
\end{longtable}

\section{Arquivos que o usuário edita}

Edite \texttt{tcc.tex}, os arquivos de \texttt{capitulos/}, \texttt{obrigatorio/}, \texttt{opcional/} e os ativos inseridos em \texttt{figuras/}. Arquivos internos de \texttt{src/} devem ser modificados apenas durante manutenção do template e acompanhados de testes nos modelos em branco e de exemplo.

Para as demais modalidades, edite a entrada escolhida e a pasta indicada no
índice dos modelos. Mantenha os arquivos de exemplos separados do trabalho real.
No TCC, os dados catalográficos continuam obrigatórios na versão final; o caminho
\texttt{opcional/ficha-catalografica.tex} é mantido por compatibilidade e não
define a obrigatoriedade do elemento.

\section{Inclusão de capítulos}

O arquivo principal usa \texttt{\textbackslash undfincludechapter}. Para ativar a seção opcional de dados ou corpus, retire o comentário da linha correspondente:

\begin{lstlisting}
\undfincludechapter{capitulos/04-dados}
\end{lstlisting}

Para adicionar outra seção primária, crie o arquivo na pasta \texttt{capitulos/} e inclua-o na posição desejada. Evite renomear arquivos internos apenas por preferência estética; caminhos inconsistentes são causa frequente de falha.

\section{Arquivos auxiliares}

LuaLaTeX e Biber criam arquivos como \texttt{.aux}, \texttt{.bcf}, \texttt{.bbl}, \texttt{.log}, \texttt{.toc}, \texttt{.lof} e \texttt{.lot}. Eles não devem ser editados manualmente. Use limpeza controlada quando houver estado inconsistente e recompile o documento completo.

\chapter{Informações Gerais do Trabalho}

As informações gerais, também chamadas de metadados, ficam no preâmbulo de \texttt{tcc.tex}. Elas alimentam páginas iniciais e propriedades do PDF.

\section{Autoria, título e vínculo}

\begin{lstlisting}
\autor{Nome do Aluno}
% \autores{Aluno 1//Aluno 2}
\titulo{Título do Trabalho}
\subtitulo{Subtítulo opcional}
\curso{Nome do curso}
\tipotrabalho{Trabalho de Conclusão de Curso}
\graupretendido{bacharel em Nome do curso}
\local{Cidade}
\ano{Ano}
\end{lstlisting}

Use \texttt{\textbackslash autor} para uma autoria. Para múltiplas autorias admitidas pelo curso, use \texttt{\textbackslash autores} e separe nomes com \texttt{//}. Não use os dois comandos ao mesmo tempo.

\section{Orientação e aprovação}

\begin{lstlisting}
\orientador{Nome do Orientador}
\coorientador{}
\areadeconcentracao{Área de concentração}
\dataaprovacao{Brasília, DF, dia de mês de ano}
\end{lstlisting}

Deixe o argumento de \texttt{\textbackslash coorientador} vazio quando não houver coorientação. A data e o formato da banca devem acompanhar a decisão institucional aplicável; os valores do modelo são exemplos editáveis.

\section{Banca examinadora}

Limpe a lista padrão e cadastre cada membro com nome, prefixo acadêmico e vínculo:

\begin{lstlisting}
\limparmembrosbanca
\membrobanca{Nome do Orientador}{Prof. Dr.}{UnDF}
\membrobanca{Nome do Membro 1}{Profa. Dra.}{Instituição}
\membrobanca{Nome do Membro 2}{Prof. Dr.}{Instituição}
\end{lstlisting}

O primeiro argumento contém o nome sem titulação. O segundo é impresso antes do nome. O terceiro aparece na linha de vínculo institucional. Cadastre apenas a composição efetivamente aprovada.

\section{Conferência}

Depois de preencher os campos, confira capa, folha de rosto, folha de aprovação e propriedades do PDF. Um erro no preâmbulo pode aparecer em mais de uma página.

\chapter{Elementos Pré-textuais no \marcaUnDFTeX}

Este capítulo descreve a implementação. Consulte a Parte I para obrigatoriedade, conteúdo e finalidade acadêmica.

\section{Montagem automática}

O arquivo principal chama, em ordem, capa, folha de rosto, verso, errata, folha de aprovação, elementos opcionais, resumos e sumário. Não duplique esses comandos dentro dos arquivos de conteúdo.

\begin{lstlisting}
\undfcover
\pretextual
\undffolhaderosto
\undfversofolhaderostoopcional
\undferrata
\undffolhadeaprovacao
\end{lstlisting}

\section{Resumos}

Preencha \texttt{obrigatorio/resumo.tex}:

\begin{lstlisting}
\undfsetcontent{resumo}{Texto do resumo.}
\undfsetcontent{palavraschave}{termo 1; termo 2; termo 3}
\end{lstlisting}

No arquivo \texttt{obrigatorio/abstract.tex}, informe título, texto e palavras-chave. Para idioma diferente do inglês, altere também o rótulo:

\begin{lstlisting}
\undfsetcontent{abstracttitle}{RESUMEN}
\undfsetcontent{abstract}{Texto em língua estrangeira.}
\undfsetcontent{keywordslabel}{Palabras clave}
\undfsetcontent{keywords}{término 1; término 2}
\end{lstlisting}

O template acrescenta o ponto ao final das palavras-chave. Por isso, informe os termos separados por ponto e vírgula, sem ponto adicional no argumento.

\section{Dados catalográficos}

O arquivo \texttt{opcional/ficha-catalografica.tex} aceita conteúdo textual, imagem ou PDF conforme os comandos públicos do modelo. Utilize apenas o material fornecido pelo fluxo institucional competente.

\begin{lstlisting}
\dadoscatalograficos{Texto fornecido pela unidade responsável.}
% \fichacatalograficaimagem{figuras/ficha.png}
% \fichacatalograficapdf{figuras/ficha.pdf}
\end{lstlisting}

\section{Errata}

\begin{lstlisting}
\errataref{Referência completa do trabalho.}
\erratatexto{%
  \erratatabela{%
    \erratalinha{16}{10}{auto-clavado}{autoclavado}
  }
}
\end{lstlisting}

\section{Elementos opcionais}

Arquivos vazios ou inexistentes não são impressos. O carregamento é feito pelo comando mostrado a seguir. Preencha apenas os elementos usados e não remova a lógica condicional do arquivo principal sem necessidade.

\begin{lstlisting}
\undfinputoptional{opcional/nome-do-arquivo}
\end{lstlisting}

\section{Listas e sumário}

\begin{lstlisting}
\listadeilustracoes
\listadefluxogramas
\listadetabelas
\undfsumario
\end{lstlisting}

O modelo em branco chama apenas o sumário por padrão. Ative listas quando forem adequadas ao trabalho e recompile até que títulos e páginas se estabilizem.

\chapter{Citações, Notas e Referências}

Este capítulo descreve a operação dos recursos bibliográficos. Os critérios para selecionar fontes, citar com integridade e elaborar referências são tratados na Parte I.

\section{Arquivo bibliográfico}

Cadastre as obras em \texttt{obrigatorio/referencias.bib} e associe o arquivo ao documento principal:

\begin{lstlisting}
\addbibresource{obrigatorio/referencias.bib}
\end{lstlisting}

Cada entrada possui uma chave única, usada pelos comandos de citação. Evite espaços, acentos e chaves duplicadas. Sempre confira no PDF se nomes, título, edição, data, endereço eletrônico e data de acesso foram processados como esperado.

\section{Citações no texto}

\begin{lstlisting}
\citarparenteses{chave}
\citarparenteses[p. 25]{chave}
\citartexto{chave}
\citartexto[p. 25]{chave}

\citardireta[p. 25]{chave}{Texto citado.}
\citardiretatexto[p. 25]{chave}{Texto citado.}
\end{lstlisting}

O argumento entre colchetes é opcional e recebe o localizador, por exemplo \texttt{p. 25}. O argumento seguinte é a chave da entrada bibliográfica.

Para citação direta longa, use:

\begin{lstlisting}
\begin{citacaolonga}[p. 25]{chave}
Texto da citação.
\end{citacaolonga}
\end{lstlisting}

Para citação de citação, prefira informar autoria e ano da original como texto,
usando uma chave apenas para a fonte efetivamente consultada:

\begin{lstlisting}
\parenciteapud{Freire}{1994}[13]{consultada}[25]
\textapud{Freire}{1994}[13]{consultada}[25]
\end{lstlisting}

Os comandos anteriores \texttt{\string\apudbib} e \texttt{\string\apudbibonline}
continuam aceitando duas chaves. Ao usar \texttt{\string\printundfreferences},
a obra original usada somente nessa citação indireta é omitida da lista; uma
citação independente da original na mesma seção bibliográfica a inclui. Essa
filtragem não é aplicada a chamadas diretas a \texttt{\string\printbibliography}.

\section{Notas de rodapé}

\begin{lstlisting}
\notarodape{Texto explicativo da nota.}
\end{lstlisting}

\section{Lista de referências}

Insira \texttt{\string\printundfreferences} após a parte textual. A lista é produzida apenas com as entradas citadas, salvo configuração bibliográfica adicional.

\section{Diagnóstico bibliográfico}

Uma citação com a chave impressa em negrito ou como marcador indefinido geralmente indica chave inexistente ou ciclo de compilação incompleto. Mensagens como \texttt{Please (re)run Biber} pedem a execução do Biber e novas passagens do LuaLaTeX.

\chapter{Figuras, Tabelas, Equações e Outros Objetos}

\section{Figura individual}

\begin{lstlisting}
\undffigura[0.75\textwidth]
  {figuras/exemplo.jpg}
  {Legenda da figura}
  {rotulo}
  [Autor, ano]
\end{lstlisting}

A largura e a fonte são opcionais. Sem indicação, o comando usa \texttt{0.8\string\textwidth} e \texttt{elaboração própria}. Referencie com \texttt{\string\figuraref\{rotulo\}}.

\section{Duas figuras independentes lado a lado}

\begin{lstlisting}
\undffigurasladoalado
  {figuras/antes.jpg}{Registro anterior}{antes}[elaboração própria]
  {figuras/depois.jpg}{Registro posterior}{depois}[elaboração própria]
\end{lstlisting}

O comando cria duas figuras independentes, cada uma com número, legenda, rótulo e fonte. Não há legenda coletiva nem subdivisões ``a'' e ``b''. A largura opcional inicial controla cada coluna.

\section{Tabela}

\begin{lstlisting}
\begin{undftabela}{Legenda da tabela}{rotulo}[Fonte]
  \begin{tabular}{ll}
    Item & Valor \\
  \end{tabular}
\end{undftabela}
\end{lstlisting}

Use \texttt{\string\tabelaref\{rotulo\}} para a chamada no texto. Tabelas extensas que atravessam páginas devem ser construídas com um ambiente multipágina apropriado, como \texttt{longtable}, pois \texttt{undftabela} é um flutuante de uma página.

\section{Gráficos e quadros}

Na versão 0.1.0, um gráfico fornecido como imagem é inserido por \texttt{\string\undffigura}; a saída o identifica na categoria geral ``Figura''. Um gráfico produzido diretamente em LaTeX pode ser colocado em um ambiente \texttt{figure}, com \texttt{\string\caption}, rótulo iniciado por \texttt{fig:} e \texttt{\string\undffonte}.

Não há ambiente público dedicado a quadro na versão 0.1.0. Não use \texttt{undftabela} apenas para obter bordas quando o objeto for semanticamente um quadro, pois a identificação sairia como ``Tabela''. Até que uma interface própria seja incorporada e testada, quadros exigem implementação manual controlada pelo mantenedor. Essa limitação está registrada para uma versão posterior e não deve ser ocultada por documentação de um comando inexistente.

\section{Equações}

\begin{lstlisting}
\equacao[energia]{E = mc^2}
\equacaosemnumero{x + y = z}
\end{lstlisting}

Somente a equação numerada e rotulada pode ser referenciada:

\begin{lstlisting}
\equacaoref{energia}
\end{lstlisting}

\section{Fluxograma}

\begin{lstlisting}
\begin{undffluxograma}{Fluxo de análise}{analise}
  \centering
  \begin{tikzpicture}
    % nós e conexões
  \end{tikzpicture}
\end{undffluxograma}
\end{lstlisting}

Até três ocorrências, o recurso integra a numeração e a lista de figuras. Acima desse total, passa a usar numeração e lista próprias. Após adicionar ou remover fluxogramas, compile novamente até os rótulos e as listas estabilizarem.

\section{Pseudocódigo}

\begin{lstlisting}
\begin{undfpseudocodigo}{MEDIA(V)}{media}{Calcular a média}
Entrada: lista V
Saida: média dos valores
...
\end{undfpseudocodigo}
\end{lstlisting}

O ambiente comum tenta iniciar o bloco apenas quando há espaço suficiente para cerca de meia página. Para um algoritmo que precise continuar na página seguinte, use \texttt{undfpseudocodigolongo}; ele permite a quebra interna do bloco. Referencie ambos com \texttt{\string\pseudocodigoref\{rotulo\}}.

\section{Apêndices, anexos e PDFs externos}

\begin{lstlisting}
\undfapendice{Título}{rotulo}
\undfanexo{Título}{rotulo}
\undfapendicepdf{Título}{arquivo.pdf}
\undfanexarpdf{Título}{arquivo.pdf}
\end{lstlisting}

Para inclusão sem criar apêndice ou anexo, use \texttt{\string\undfincluirpdf\{arquivo.pdf\}}. Os comandos de referência são \texttt{\string\apendiceref} e \texttt{\string\anexoref}.

\section{Imagens e paginação}

Os comandos de figura verificam o tamanho do arquivo e emitem avisos para imagens pesadas. Uma figura fixada que não caiba no espaço restante avança inteira para a próxima página, podendo deixar área em branco. Isso preserva a unidade do objeto. Ajuste a largura, reposicione a chamada no texto ou use uma imagem mais compacta; não force sobreposição com texto ou margens.

\chapter{Compilação e Geração do PDF}

\section{Fluxo automatizado}

Na raiz do projeto, execute:

\begin{lstlisting}
latexmk -outdir=build/perfis
\end{lstlisting}

O arquivo \texttt{.latexmkrc} seleciona LuaLaTeX, aciona o Biber quando necessário e define \texttt{tcc.tex} como entrada padrão. Para o exemplo completo:

\begin{lstlisting}
latexmk -outdir=build/perfis tcc-exemplo-undf.tex
\end{lstlisting}

\section{Fluxo manual}

Quando o \texttt{latexmk} não estiver disponível, use:

\begin{lstlisting}
New-Item -ItemType Directory -Force build/perfis | Out-Null
lualatex -output-directory=build/perfis tcc.tex
biber --input-directory=build/perfis --output-directory=build/perfis tcc
lualatex -output-directory=build/perfis tcc.tex
lualatex -output-directory=build/perfis tcc.tex
\end{lstlisting}

No PowerShell, as saídas ficam em \texttt{build/perfis/}. O script
\texttt{scripts/verificar-perfis.ps1} automatiza esse ciclo no Windows, inclusive
sem Perl. O Biber deve receber o nome-base, sem a extensão \texttt{.tex}.
Passagens adicionais podem ser necessárias após alterações em sumário, listas,
citações ou referências cruzadas.

\section{Overleaf}

Envie a estrutura completa do projeto e selecione LuaLaTeX como compilador. O arquivo \texttt{.latexmkrc} acompanha o repositório. Confirme no painel de logs se a bibliografia foi processada e se não restaram referências indefinidas.

\section{Arquivos auxiliares}

Arquivos como \texttt{.aux}, \texttt{.bcf}, \texttt{.bbl}, \texttt{.toc}, \texttt{.lof}, \texttt{.lot} e \texttt{.lofl} registram dados entre as passagens. Eles não fazem parte do texto autoral e podem ser regenerados. Para limpar artefatos pelo \texttt{latexmk}, use \texttt{latexmk -outdir=build/perfis -c}; faça nova compilação completa em seguida.

\section{Critério de conclusão}

Considere a geração tecnicamente estabilizada somente quando: o comando terminar sem erro fatal; o PDF abrir normalmente; não houver pedido de nova execução do Biber; sumário, listas e referências cruzadas estiverem resolvidos; e a inspeção visual das páginas críticas tiver sido feita.

\chapter{Validações e Avisos}

\section{Estimativa de ocupação dos elementos pré-textuais}

O arquivo principal chama, antes da capa:

\begin{lstlisting}
\undfvalidateonepagecontents
\end{lstlisting}

O comando estima o espaço ocupado por resumo, resumo em língua estrangeira, dedicatória, agradecimentos, epígrafe, errata e ficha catalográfica textual. A medição é preventiva: caixas, fontes, imagens e outros elementos podem fazer o resultado visual diferir da estimativa. Sempre confira o PDF.

A estimativa é um recurso técnico do projeto e não cria, por si só, obrigação normativa ou institucional de limitar o elemento a uma página. Quando houver limite definido pelo curso ou por decisão institucional aplicável, essa origem deve ser identificada separadamente.

\section{Resumo e resumo em língua estrangeira}

Por padrão, excesso estimado de conteúdo gera aviso técnico no log e a compilação continua. Para transformar esse caso em erro intencional durante uma revisão controlada, use antes da validação:

\begin{lstlisting}
\bloquearresumoslongos
\undfvalidateonepagecontents
\end{lstlisting}

Para restabelecer explicitamente o modo de aviso:

\begin{lstlisting}
\avisarresumoslongos
\end{lstlisting}

O sistema também estima a contagem de palavras e avisa quando o resumo parece ficar fora da faixa configurada. Essa contagem pode divergir de editores de texto diante de comandos, siglas, números e palavras compostas.

O modo de erro é uma ferramenta optativa de revisão e não representa o comportamento oficial de entrega. Uma recomendação normativa ou editorial, isoladamente, não justifica bloquear a compilação.

\section{Conteúdo ausente}

\texttt{\string\undfwarnifempty} produz aviso quando um conteúdo obrigatório não foi preenchido. Arquivos opcionais vazios ou inexistentes são ignorados por \texttt{\string\undfinputoptional}.

\section{Imagens}

Os comandos de figura podem avisar sobre arquivos acima dos limiares configurados. Para suspender e reativar a verificação:

\begin{lstlisting}
\undfdesabilitaravisosimagem
\undfhabilitaravisosimagem
\end{lstlisting}

Esses avisos avaliam peso do arquivo, não nitidez, contraste, legibilidade nem adequação de resolução. A conferência visual continua obrigatória.

\section{Natureza dos diagnósticos}

\begin{description}
  \item[Erro] interrompe ou invalida a geração; deve ser corrigido.
  \item[Aviso] sinaliza risco ou pendência; exige avaliação humana.
  \item[Informação] registra uma decisão do processamento e não implica defeito.
  \item[Estimativa] aproxima ocupação ou contagem e sempre exige conferência humana.
\end{description}

Requisito normativo, decisão institucional validada, recomendação e regra de curso pertencem à documentação acadêmica; aviso, estimativa e erro de compilação pertencem ao diagnóstico técnico. Não trate um PDF gerado como prova automática de conformidade. O template automatiza forma e diagnósticos limitados, mas não avalia conteúdo, método, autoria ou decisão institucional.

\chapter{Solução de Problemas}

\begin{longtable}{>{\raggedright\arraybackslash}p{0.25\textwidth} >{\raggedright\arraybackslash}p{0.31\textwidth} >{\raggedright\arraybackslash}p{0.34\textwidth}}
\toprule
\textbf{Sintoma} & \textbf{Causa provável} & \textbf{Ação recomendada} \\
\midrule
\endhead
Classe não encontrada & Compilação fora da raiz ou pacote incompleto & Compile na raiz e confira \texttt{undftex.cls} e a pasta \texttt{src/}. \\
\addlinespace
LuaLaTeX exigido & Motor incorreto & Selecione LuaLaTeX no editor ou serviço. \\
\addlinespace
Citações indefinidas & Chave ausente ou Biber não executado & Confira o \texttt{.bib}, execute Biber e repita LuaLaTeX. \\
\addlinespace
Referência cruzada com ``??'' & Passagens insuficientes ou rótulo divergente & Confira o rótulo e recompile até estabilizar. \\
\addlinespace
Sumário ou listas desatualizados & Arquivos auxiliares antigos & Faça novas passagens; se persistir, limpe auxiliares e recompile. \\
\addlinespace
Figura não encontrada & Caminho, extensão ou maiúsculas divergentes & Confira o nome real e use caminho relativo à raiz. \\
\addlinespace
Fonte ausente ou substituída & Distribuição incompleta ou cache de fontes & Atualize a distribuição, confira o log e não aceite substituição silenciosa sem revisar o PDF. \\
\addlinespace
Caractere especial causa erro & Símbolo reservado ou arquivo fora de UTF-8 & Preserve UTF-8 e use a sintaxe LaTeX adequada para caracteres reservados. \\
\addlinespace
Grande espaço antes de figura & Objeto fixado não cabe no restante da página & Reduza a largura, reposicione a chamada ou aceite o avanço integral. \\
\addlinespace
Página em branco inesperada & Quebra explícita, inclusão vazia ou ajuste de objeto & Localize a transição no fonte e verifique \texttt{clearpage}, arquivos opcionais e objetos anteriores. \\
\addlinespace
Conteúdo fora da margem & Linha de código, URL, tabela ou imagem larga & Reduza ou permita quebra; procure \texttt{Overfull} no log e confira a página. \\
\addlinespace
Pseudocódigo ocupa página seguinte & Bloco não cabe ou é longo & Use \texttt{undfpseudocodigolongo} quando a quebra interna for necessária. \\
\addlinespace
Fluxograma com lista incorreta & Contagem mudou desde a última passagem & Recompile até desaparecer o aviso de contagem. \\
\addlinespace
Resumo gera aviso de uma página & Estimativa de altura excedida & Revise o texto e confira o PDF; o padrão não bloqueia a compilação. \\
\addlinespace
PDF externo com página cortada & Dimensões ou opções de inclusão incompatíveis & Ajuste opções de \texttt{pdfpages} no argumento opcional e confira todas as páginas. \\
\addlinespace
Compilação interrompida & Erro de sintaxe, arquivo ausente ou ferramenta indisponível & Corrija a primeira mensagem iniciada por \texttt{!} e execute novamente o ciclo completo. \\
\bottomrule
\end{longtable}

\section{Leitura do log}

Procure primeiro a primeira linha iniciada por \texttt{!}, pois erros posteriores podem ser consequência dela. Em seguida, busque \texttt{Undefined}, \texttt{Warning}, \texttt{Overfull} e pedidos de nova compilação. Registre o comando executado e a primeira mensagem relevante ao solicitar suporte.

\section{Isolamento do problema}

Quando a causa não estiver clara, reproduza o erro no menor conjunto possível: comente temporariamente a inclusão do capítulo suspeito, reduza o conteúdo até identificar o trecho e restaure as partes gradualmente. Faça esse diagnóstico em cópia ou sob controle de versão para não perder conteúdo.

\chapter{Checklist Técnico de Fechamento}

\section{Arquivos e versão}

\begin{itemize}[label=$\square$]
  \item Confirme que o trabalho usa uma cópia íntegra do UnDF-TeX 0.1.0.
  \item Confira os arquivos de entrada, capítulos, imagens e bibliografia.
  \item Remova apenas artefatos temporários; preserve fontes e arquivos necessários à reprodução.
  \item Verifique se não há dados de exemplo remanescentes.
\end{itemize}

\section{Compilação}

\begin{itemize}[label=$\square$]
  \item Compile com LuaLaTeX e Biber pelo fluxo recomendado.
  \item Repita as passagens até sumário, listas e rótulos estabilizarem.
  \item Corrija todos os erros fatais e avalie cada aviso relevante.
  \item Confirme que não há citações ou referências cruzadas indefinidas.
\end{itemize}

\section{Conferência visual}

\begin{itemize}[label=$\square$]
  \item Abra o PDF gerado, não apenas a visualização do editor.
  \item Confira capa, folha de rosto, aprovação, resumos e sumário.
  \item Examine transições de página, objetos extensos e páginas com pouco conteúdo.
  \item Confira nitidez, legenda, fonte e chamada de cada objeto.
  \item Verifique links, marcadores, numeração e páginas incluídas externamente.
  \item Confirme ordem de leitura, contraste, legibilidade e alternativas textuais quando o fluxo de produção as suportar.
  \item Faça uma leitura em escala de 100\% e outra em visão geral de miniaturas.
\end{itemize}

\section{Entrega técnica}

\begin{itemize}[label=$\square$]
  \item Gere o PDF a partir da versão efetivamente aprovada para entrega.
  \item Use nome de arquivo claro, sem marcas como \texttt{final-final}.
  \item Confirme abertura em outro leitor de PDF.
  \item Arquive o código-fonte, a bibliografia e os recursos usados na geração.
\end{itemize}

\makeatletter
\gdef\@chapapp{Capítulo}
\makeatother
\chapter{Compatibilidade, Licença e Manutenção}

\section{Matriz de compatibilidade}

\begin{longtable}{p{0.27\textwidth}p{0.25\textwidth}p{0.38\textwidth}}
\toprule
\textbf{Documento ou componente} & \textbf{Versão} & \textbf{Situação} \\
\midrule
\endhead
Guia de Uso do UnDF-TeX & 1.0.0 & Compatível com a versão técnica indicada nesta edição. \\
Manual de TCC da UnDF & 1.0.0 & Documento acadêmico autônomo relacionado. \\
UnDF-TeX & 0.1.0 & Versão de referência deste guia. \\
LuaLaTeX e Biber & Distribuições atuais compatíveis & Devem oferecer os pacotes carregados pelo projeto. \\
\bottomrule
\end{longtable}

\section{Versionamento}

Mudanças na classe ou nos pacotes podem alterar comandos e aparência. Antes de atualizar um trabalho em andamento, registre a versão atual, crie uma cópia ou uma ramificação e compare os PDFs antes e depois. Não misture arquivos internos de versões diferentes.

Consulte o \href{https://github.com/igormrqs/undftex-publico/blob/main/CHANGELOG.md}{histórico público de alterações}
e a \href{https://github.com/igormrqs/undftex-publico}{página do repositório}
para identificar a origem do pacote. As versões editoriais 1.0.0 do manual
e do guia não representam, por si só, atualização do template para 1.0.0.

\section{Licença}

A LPPL 1.3c ou posterior aplica-se aos arquivos de código expressamente identificados. Consulte \texttt{LICENSE} e \texttt{LICENSE-LPPL-1.3c} antes de modificar ou redistribuir esses arquivos.

Marcas, logotipos, elementos de identidade institucional, fontes, imagens, documentos oficiais e recursos de terceiros permanecem sujeitos às respectivas regras de uso e não são automaticamente licenciados pela LPPL. A documentação textual somente deve ser redistribuída sob a licença que vier a ser expressamente indicada para ela.

\section{Responsabilidades de manutenção}

Uma alteração técnica deve atualizar código, exemplo, testes pertinentes, documentação de uso, histórico e indicação de compatibilidade. A validação deve incluir compilação do modelo em branco e do exemplo completo, comparação visual e revisão dos avisos do log.

\section{Relato de problema}

Informe versão, ambiente, entrada, comando, primeira mensagem e trecho mínimo. Em defeitos visuais, inclua página ou captura; omita dados pessoais e conteúdo sigiloso.

Problemas relativos à base pública podem ser registrados na
\href{https://github.com/igormrqs/undftex-publico/issues}{área de issues do repositório}.
Para recursos experimentais, use o canal de acompanhamento disponibilizado
pelo responsável pelo desenvolvimento privado.

\chapter{Referência Rápida de Comandos}

As tabelas reúnem a interface pública mais usada. Nos exemplos, argumentos entre colchetes são opcionais; argumentos entre chaves são obrigatórios.

\section{Informações gerais}

\begin{longtable}{>{\ttfamily\footnotesize\raggedright\arraybackslash}p{0.40\textwidth}p{0.20\textwidth}p{0.30\textwidth}}
\toprule
\normalfont\textbf{Comando} & \textbf{Finalidade} & \textbf{Uso e observações} \\
\midrule
\endhead
\string\autor & Autor único & Parâmetro: nome. Use \string\autores para trabalho coletivo. \\
\string\autores & Dois ou mais autores & Parâmetro: nomes separados por \texttt{//}. \\
\string\titulo & Título principal & Parâmetro: texto sem a pontuação do subtítulo. \\
\string\subtitulo & Subtítulo opcional & Parâmetro: texto; deixe vazio quando inexistente. \\
\string\curso & Curso & Parâmetro: nome por extenso. \\
\string\tipotrabalho & Tipo documental & Parâmetro: tipo do trabalho. \\
\string\graupretendido & Grau pretendido & Parâmetro: texto que completa a natureza. \\
\string\local & Local & Parâmetro: cidade. \\
\string\ano & Ano & Parâmetro: quatro algarismos. \\
\string\orientador & Orientador & Parâmetro: nome sem prefixo acadêmico. \\
\string\coorientador & Coorientador & Parâmetro: nome; pode ficar vazio. \\
\string\areadeconcentracao & Área de concentração & Parâmetro: área, quando pertinente. \\
\string\dataaprovacao & Data de aprovação & Parâmetro: local e data por extenso. \\
\string\membrobanca & Membro da banca & Parâmetros: nome, prefixo e instituição. \\
\bottomrule
\end{longtable}

\section{Conteúdo e arquivos}

\begin{longtable}{>{\ttfamily\footnotesize\raggedright\arraybackslash}p{0.40\textwidth}p{0.20\textwidth}p{0.30\textwidth}}
\toprule
\normalfont\textbf{Comando} & \textbf{Finalidade} & \textbf{Uso e observações} \\
\midrule
\endhead
\string\undfsetcontent & Registrar conteúdo & Parâmetros: chave e conteúdo. \\
\string\undfusecontent & Recuperar conteúdo & Parâmetro: chave. \\
\string\undfifcontent & Testar conteúdo & Parâmetros: chave, saída preenchida e saída vazia. \\
\string\undfinputoptional & Incluir arquivo opcional & Parâmetro: caminho sem \texttt{.tex}. \\
\string\undfincludechapter & Incluir capítulo & Parâmetro: caminho sem \texttt{.tex}. \\
\string\undfwarnifempty & Avisar ausência & Parâmetros: chave e mensagem; não interrompe. \\
\string\undfincluirpdf & Incluir PDF integral & Parâmetro: arquivo; aceita opções em colchetes. \\
\bottomrule
\end{longtable}

\section{Elementos pré e pós-textuais}

\begin{longtable}{>{\ttfamily\footnotesize\raggedright\arraybackslash}p{0.40\textwidth}p{0.22\textwidth}p{0.28\textwidth}}
\toprule
\normalfont\textbf{Comando} & \textbf{Resultado} & \textbf{Observações} \\
\midrule
\endhead
\string\undfcover & Capa & Usa as informações gerais. \\
\string\undffolhaderosto & Folha de rosto & Usa natureza e orientação. \\
\string\undfversofolhaderostoopcional & Ficha no verso & Só aparece com conteúdo configurado. \\
\string\undferrata & Errata & Só aparece quando preenchida. \\
\string\undffolhadeaprovacao & Folha de aprovação & Usa data e membros da banca. \\
\string\listadeilustracoes & Lista de figuras & Gerada a partir das legendas. \\
\string\listadefluxogramas & Lista própria & Exibida quando a contagem exige. \\
\string\listadetabelas & Lista de tabelas & Gerada a partir das legendas. \\
\string\undfsumario & Sumário & Exige novas passagens após mudanças. \\
\string\printundfreferences & Referências & Usa a base ligada por \string\addbibresource. \\
\string\undfapendice & Apêndice & Parâmetros: título e rótulo. \\
\string\undfanexo & Anexo & Parâmetros: título e rótulo. \\
\string\undfapendicepdf & Apêndice em PDF & Parâmetros: título e arquivo. \\
\string\undfanexarpdf & Anexo em PDF & Parâmetros: título e arquivo. \\
\bottomrule
\end{longtable}

\section{Citações e referências}

\begin{longtable}{>{\ttfamily\footnotesize\raggedright\arraybackslash}p{0.40\textwidth}p{0.22\textwidth}p{0.28\textwidth}}
\toprule
\normalfont\textbf{Comando} & \textbf{Resultado} & \textbf{Erro comum} \\
\midrule
\endhead
\string\citarparenteses & Citação parentética & Localizador opcional e chave obrigatória. \\
\string\citartexto & Autor integrado ao texto & Localizador opcional e chave obrigatória. \\
\string\citardireta & Citação curta parentética & Localizador, chave e texto. \\
\string\citardiretatexto & Citação curta narrativa & Localizador, chave e texto. \\
\string\citacaolonga & Bloco de citação longa & Localizador opcional e chave obrigatória. \\
\string\parenciteapud & Citação de citação & Autoria/ano originais e chave da consultada. \\
\string\textapud & Autoria integrada ao texto & Mesmos dados e localizadores opcionais. \\
\string\apudbib & Interface compatível anterior & Dois localizadores e duas chaves. \\
\string\fullref & Referência completa pontual & Parâmetro: chave. \\
\string\notarodape & Nota de rodapé & Parâmetro: texto. \\
\bottomrule
\end{longtable}

\section{Objetos e referências cruzadas}

\begin{longtable}{>{\ttfamily\footnotesize\raggedright\arraybackslash}p{0.43\textwidth}p{0.20\textwidth}p{0.27\textwidth}}
\toprule
\normalfont\textbf{Comando ou ambiente} & \textbf{Resultado} & \textbf{Observações} \\
\midrule
\endhead
\string\undfsecao & Seção rotulada & Parâmetros: título e rótulo. \\
\string\undfsubsecao & Subseção rotulada & Parâmetros: título e rótulo. \\
\string\equacao & Equação numerada & Rótulo opcional e expressão. \\
\string\undffigura & Figura & Largura, arquivo, legenda, rótulo e fonte. \\
\string\undffigurasladoalado... & Duas figuras independentes & Cada figura recebe seus próprios dados. \\
\string\undftabela & Tabela flutuante & Legenda, rótulo e fonte opcional. \\
\string\undffluxograma & Fluxograma & Legenda, rótulo e fonte opcional. \\
\string\undfpseudocodigo & Pseudocódigo comum & Assinatura, rótulo e título. \\
\string\undfpseudocodigolongo & Pseudocódigo multipágina & Mesmos parâmetros; permite quebra. \\
\string\figuraref & Chamada de figura & Parâmetro: rótulo sem \texttt{fig:}. \\
\string\tabelaref & Chamada de tabela & Parâmetro: rótulo sem \texttt{tab:}. \\
\string\equacaoref & Chamada de equação & Parâmetro: rótulo sem \texttt{eq:}. \\
\string\pseudocodigoref & Chamada de pseudocódigo & Parâmetro: rótulo sem \texttt{alg:}. \\
\bottomrule
\end{longtable}

\section{Validações}

\begin{longtable}{>{\ttfamily\footnotesize\raggedright\arraybackslash}p{0.40\textwidth}p{0.24\textwidth}p{0.26\textwidth}}
\toprule
\normalfont\textbf{Comando} & \textbf{Efeito} & \textbf{Observação} \\
\midrule
\endhead
\string\undfvalidateonepagecontents & Executa estimativas & Chame antes da capa. \\
\string\bloquearresumoslongos & Excesso vira erro & Opção de revisão, não padrão. \\
\string\avisarresumoslongos & Excesso permanece aviso & Comportamento padrão. \\
\string\undfdesabilitaravisosimagem & Suspende avisos de peso & Não altera a imagem. \\
\string\undfhabilitaravisosimagem & Reativa avisos & Ativo por padrão. \\
\bottomrule
\end{longtable}

\section{Exemplos e resultados esperados}

As duas tabelas seguintes se complementam. A primeira apresenta finalidade, parâmetros e uso básico; a segunda registra resultado esperado, erro comum e solução. Os comandos de apoio menos frequentes permanecem documentados nas tabelas temáticas anteriores.

\begin{landscape}
\begin{longtable}{>{\ttfamily\footnotesize\raggedright\arraybackslash}p{0.25\linewidth}p{0.19\linewidth}p{0.22\linewidth}>{\ttfamily\footnotesize\raggedright\arraybackslash}p{0.26\linewidth}}
\toprule
\normalfont\textbf{Comando} & \textbf{Finalidade} & \textbf{Parâmetros} & \normalfont\textbf{Uso básico} \\
\midrule
\endhead
\string\autor & Registrar autoria & Nome obrigatório & \string\autor\{Ana Lima\} \\
\string\citarparenteses & Inserir citação parentética & Localizador opcional e chave obrigatória & \string\citarparenteses\{silva2025\} \\
\string\undffigura & Inserir figura & Largura opcional, arquivo, título, rótulo e fonte opcional & \string\undffigura\{img\}\{Título\}\{id\} \\
\string\equacao & Inserir equação & Rótulo opcional e expressão obrigatória & \string\equacao[e]\{E=mc\string^2\} \\
\string\undfsecao & Criar seção rotulada & Título e rótulo obrigatórios & \string\undfsecao\{Objetivos\}\{obj\} \\
\string\bloquearresumoslongos & Ativar erro optativo de revisão & Nenhum & \string\bloquearresumoslongos \\
\bottomrule
\end{longtable}

\begin{longtable}{>{\ttfamily\footnotesize\raggedright\arraybackslash}p{0.25\linewidth}p{0.23\linewidth}p{0.21\linewidth}p{0.23\linewidth}}
\toprule
\normalfont\textbf{Comando} & \textbf{Resultado esperado} & \textbf{Erro comum} & \textbf{Solução} \\
\midrule
\endhead
\string\autor & Autoria nas páginas iniciais & Deixar o exemplo padrão & Substituir pelo nome real antes de compilar \\
\string\citarparenteses & Chamada autor-data vinculada à bibliografia & Usar chave inexistente & Conferir a chave no arquivo bibliográfico e executar o Biber \\
\string\undffigura & Figura numerada, com legenda, rótulo e fonte & Informar caminho incorreto & Conferir nome, extensão e caminho relativo do arquivo \\
\string\equacao & Equação numerada e referenciável & Repetir o rótulo & Usar um identificador único e recompilar \\
\string\undfsecao & Seção numerada e referenciável & Repetir o rótulo & Adotar rótulo único e descritivo \\
\string\bloquearresumoslongos & Excesso estimado tratado como erro & Ativar como comportamento comum de escrita & Reservar o comando para revisão controlada; o padrão oficial é aviso \\
\bottomrule
\end{longtable}
\end{landscape}


\chapter*{Referências Técnicas e Documentos Relacionados}
\addcontentsline{toc}{chapter}{Referências Técnicas e Documentos Relacionados}

\begin{itemize}
  \item \emph{Manual de Trabalhos de Conclusão de Curso da UnDF}, versão 1.0.0;
  \item documentação pública do projeto: \href{https://github.com/igormrqs/undftex-publico/blob/main/README.md}{README}, \href{https://github.com/igormrqs/undftex-publico/blob/main/docs/guia-de-uso.md}{guia de uso} e \href{https://github.com/igormrqs/undftex-publico/blob/main/docs/estrutura-do-projeto.md}{estrutura do projeto};
  \item documentação do LuaLaTeX, Biber, BibLaTeX e dos pacotes carregados pela classe;
  \item LaTeX Project Public License, versão 1.3c ou posterior;
  \item código-fonte da base estável do UnDF-TeX: \url{https://github.com/igormrqs/undftex-publico}.
\end{itemize}

\chapter*{Controle Documental}
\addcontentsline{toc}{chapter}{Controle Documental}

\noindent\begin{tabularx}{\textwidth}{@{}>{\bfseries}p{0.31\textwidth}Y@{}}
\toprule
Título & Manual de Trabalhos de Conclusão de Curso da UnDF e Guia de Uso do UnDF-TeX \\
Versão do manual & 1.0.0 \\
Compatibilidade & UnDF-TeX 1.0.0 \\
Situação & Minuta para validação institucional \\
Data & 20 de outubro de 2026 \\
Repositório público & \url{https://github.com/igormrqs/undftex-publico} \\
Escopo & Orientações acadêmicas, metodológicas, normativas e técnicas \\
\bottomrule
\end{tabularx}

\end{document}
